%% 波-9 音の大きさをはかる：音速・デシベル・騒音（第24回編集委員会 物理第3巻：組 波W5）。★1〜4、3段にしない。
%% 記号（計画書 C-1・B の波-9。この章の中だけの分け方）：
%%   音の速さ（波の速さ）c（熱-4 ⑥は a と書いた。この巻は c：W-3）。セルシウス温度 θ〔℃〕、絶対温度 T〔K〕。
%%   この章では周期 T を書かず、振動数 f で書く（T は絶対温度だけ）。比熱比 γ・比気体定数 R_s は熱-3 のとおり（空気 1.4・287 J/(kg·K)）。
%%   音圧 p（大気圧からのずれ。ことわらないかぎり実効値）、大気圧 p_0、空気の密度 ρ_0、空気が動く速さ u、空気のずれ ξ。
%%   体積弾性率 K（断熱は K_s：熱-3 ⑧）、ヤング率 E、横弾性係数 G（力-8・力-11）。
%%   音の強さ I〔W/m²〕、音響パワー W〔W〕。音圧レベル L_p、音の強さのレベル L_I、音響パワーレベル L_W〔dB〕。
%%   距離 R（波-7 ②と同じ。修3-4 で r から直した：P4-10）、断面積・表面積 S、体積 V、面密度 m〔kg/m²〕、透過損失 TL（立体の2文字）、吸音率 α、吸音力 A〔m²〕（この章は振幅 A を使わない）、
%%   残響時間 t_R（専門書の T・T_60。絶対温度 T とぶつけない）。
%% 公式の持ち主（計画書 A）：A1・A4（波-5 ①と分担）・A17・A47・A62 はこの章。気体の音速 A55 は熱-4 ⑥（W-7）。
%%   √(E/ρ)・√(K/ρ)（A54）は波-5 ⑤。I＝p²/(ρc) と Z＝ρc は波-7（ここでは「押す力×速さ」で短く示して使う）。
%%   ニュートンとラプラスの話は熱-4 ⑥のコラムにあるので、くり返さない（指すだけ）。
%%   修3-4（第25回 X-5）：流体の波動方程式と音速の導き方は波-5 ⑤、点の音源の逆2乗の法則は波-7 ②が持ち主。
%%     ③は導き直し（管の中の層の図・運動方程式）をけずって波-5 ⑤を指し、⑦は波-7 ②を指す（音響パワーは W と書くと断る）。
%%     公式の欄に「音の速さは約 340 m/s（中学理科）」（計画書 A5。ラベルは⑤の読み方の中）を足した。
%% ほかの巻との約束：
%%   熱-4 ⑥・例題8：空気 20 ℃で 343 m/s、ヘリウム 1008 m/s（U-10）、「1 K で約 0.6 m/s（詳しくは波-9）」→ ⑤で受ける。
%%   熱-3 ⑧：K_s＝γp（空気 0.14 MPa）→ ③。熱-14 ③：窒素の平均の速さ 471 m/s、音の速さはその約0.73倍 → ⑤のせいかくに言うと。
%%   数学 3-12⑦：ささやき 30 dB・図書館 40 dB・会話 60 dB・工場 90 dB・ガード下 100 dB、60＋60→63 dB → ⑥⑦・例題8・9。
%%   力-8・力-12：水の体積弾性率 約2.2 GPa、鋼 E＝205 GPa・ρ＝7850 kg/m³・K 約170 GPa、力-11：G＝79 GPa。力-16：鋼の棒 約5100 m/s。
%% 数値はすべて Python で確かめた：checks/物理第3巻/波-9/calc_p3_09.py。
%% 出典（第24回の追記：書くときに確かめる。11月5日、組 波W5。この作業環境では WebFetch が DNS で届かなかった
%%   ので、curl で読めたもの（環境省・厚生労働省）は本文で、ほかは WebSearch で一次資料を引く資料を2つ以上照らし合わせた。
%%   JIS の閲覧（webdesk.jsa.or.jp）のプレビューは curl では HTML しか返らず、表は読めなかった）：
%%   ・騒音に係る環境基準：環境庁告示第64号（平成10年9月30日、平成24年3月30日 環告54 まで改正）。curl で環境省 www.env.go.jp/kijun/oto1-1.html の本文を読んで確かめた（評価は等価騒音レベル・A 特性・JIS Z 8731）。AA 昼50・夜40、A及びB 昼55・夜45、C 昼60・夜50 dB 以下、
%%     幹線交通を担う道路に近接する空間 昼70・夜65 dB 以下。照らし合わせ：環境省 www.env.go.jp/kijun/oto1-1.html（検索結果の抜粋）、
%%     名古屋市 62Dhoho_shiryo.pdf、仙台市 7soon.html、三重県 001117616.pdf。〔確〕
%%   ・騒音障害防止のためのガイドライン：令和5年4月20日 基発0420第2号で改訂。等価騒音レベル 85 dB 以上になる可能性の大きい作業場（別表）。
%%     厚生労働省のパンフレット www.mhlw.go.jp/content/001195349.pdf を curl で読んで確かめた（第Ⅰ管理区分 85 dB 未満、第Ⅱ 85〜90 dB）。〔確〕
%%   ・A 特性：IEC 61672-1:2013（JIS C 1509-1:2017）の表3 の値を、python-acoustics の data/iec_61672_1_2013.csv（表3 を写したもの。
%%     curl で読んだ）で確かめた：63 Hz −26.2、100 Hz −19.1、1 kHz 0、4 kHz ＋1.0 dB。本文の計算は規格の附属書の式（20.6・107.7・737.9・12194 Hz、
%%     −2.00 dB）で、表の値と小数第1位まで合う（12.5 Hz・16 Hz は正確な中心周波数を使う表と0.3 dB ちがう）。〔確〕
%%   ・基準音圧 20 μPa（1 kHz の最小可聴音圧）・基準の強さ 10^-12 W/m²、120 dB≈20 Pa・130 dB≈63 Pa：MIT OCW 4.401 lec15、宮崎大学の講義資料。〔確〕
%%   ・聞こえる範囲 約20 Hz〜20 kHz、20 kHz 以上を超音波：医学書院の用語解説、TOA「音の豆知識」2-4。〔確〕
%%   ・フォンとソーン（40 フォン＝1 ソーン、10 フォンごとに約2倍）：S. S. Stevens, Psychological Review 43, 405–416 (1936)（定義のもと）、
%%     HyperPhysics「Phon」、J. D. Cook の換算の解説（20〜120 フォンで実験で確かめられた）。〔確〕
%%   ・外有毛細胞の伸び縮み（Brownell ら 1985年）、耳音響放射（Kemp 1978年）、新生児の聞こえの検査：Acoustics Today 2017 Spring、
%%     Canadian Audiologist「Who Measures Otoacoustic Emissions?」。増幅のしくみの全体は確定していない（AudiologyOnline の総説）。〔確〕
%%   ・水の音の速さ：Marczak（1997年）の純水の式で 20 ℃ 1482 m/s・8 ℃ 1439 m/s（checks で計算）。20 ℃ 約1480 m/s、海水 約1500 m/s は
%%     DOSITS「How fast does sound travel?」ほか。〔確〕
%%   ・コラドンとシュトゥルム（1826年、ジュネーブ湖、水温 8 ℃、1435 m/s、鐘と火薬）：DOSITS「The First Studies of Underwater Acoustics」、
%%     arXiv:1710.03777。船の間の距離は資料で 10 km・約16 km などとちがうので、本文は「10 km あまり」。〔確〕
%%   ・セイビン（1895年 フォッグ美術館の講堂、サンダース劇場の座ぶとん、1900年 ボストン・シンフォニーホール）：
%%     APS News 2011年1月「This Month in Physics History」、Harvard Physics「2023-acoustics.pdf」。〔確〕
%%   ・教室の残響時間 0.6 s 以下（ANSI S12.60-2002、容積 283 m³ 未満の教室）：Armstrong・Larson Davis の解説。コンサートホール 約2 s 前後
%%     （1.4〜2.6 s の範囲：Canadian Acoustics の論文、Commercial Acoustics の目標値の表）。〔確〕
%%   ・コンクリートの吸音率（500 Hz）0.02〜0.04：Commercial Acoustics の吸音率の表（塗装 0.02、無塗装 0.04）。〔確〕
%%   ・板ガラスのコインシデンス：日建連の音の用語集 ka-0600・ka-0400（限界周波数より上で質量則より遮音が下がる）。5 mm で約2400 Hz は
%%     E＝72 GPa・ν＝0.23・ρ＝2500 kg/m³ から計算。ガラスの密度 約2500 kg/m³（ソーダ石灰ガラスの一般的な値）。〔確（計算）〕
%%   ・鋼の縦波 約5900 m/s：三菱重工の超音波探傷の資料、Olympus の音速表（1020 鋼 5890 m/s）。横波 約3200 m/s は E・ν から計算（3170 m/s）。〔確〕
%%   ・暗騒音の扱い（差が10 dB 以上なら補正しない）：JIS Z 8731 の本文は確かめられなかった。本文は表2 の計算（0.4 dB 以下）を理由に書いた。
%%   〔要確認〕のまま：なし（Wikipedia は出典にしていない）。
\makeatletter
\tikzset{rxj lab/.style={font=\sffamily\footnotesize, inner sep=1pt},
  rxj arr/.style={-{Stealth[length=3.6pt, width=3.2pt]}, line width=1pt},
  rxj big/.style={-{Stealth[length=4.2pt, width=3.8pt]}, line width=1.4pt},
  rxj dim/.style={{Stealth[length=3pt]}-{Stealth[length=3pt]}, line width=0.5pt, bkGray}}
% 空気の粒を、密・疎で並べる  \rxjair{左はし x}{右はし x}{中心の y}{高さの半分}{振幅（mm）}{波長（mm）}
\newcommand\rxjair[6]{%
  \foreach \rxji in {0,...,60} {%
    \pgfmathsetmacro\rxjxx{#1+\rxji*(#2-#1)/60}%
    \pgfmathsetmacro\rxjdx{-#5*sin(360*((\rxjxx-#1)/#6))}%
    \pgfmathsetmacro\rxjpx{\rxjxx+\rxjdx}%
    \foreach \rxjj in {-2,...,2} {%
      \pgfmathsetmacro\rxjyy{#3+\rxjj*#4/5+0.6*mod(\rxji,2)}%
      \fill[bkBlue!70!black] (\rxjpx,\rxjyy) circle (0.55);}}}
\makeatother

\phchapter{9}{音の大きさをはかる：\\音速・デシベル・騒音}{雷が光ってから、なぜ少しおくれて音が聞こえる？}{\bklevelrange{1}{4}}{音の大きさをはかる：音速・デシベル・騒音}
\bkmichishirube{3-12,5-17,6-4,熱-3,熱-4,熱-14,波-5,波-7}{波-10,波-11}

\begin{goal}
\begin{itemize}[label={\color{bkTheme}●}, leftmargin=1.6\zw, itemsep=0pt, topsep=1pt]
\item \bklevel{1}\enspace 雷やこだまの時間から、音の速さを使って距離を求められる。
\item \bklevel{3}\enspace 音の大きさを dB で表し、音を合わせたときや、離れたときの大きさを計算できる。
\item \bklevel{4}\enspace 縮みにくさと密度から気体・液体・固体の音の速さを求め、遮音と残響を見積もれる。
\end{itemize}
\end{goal}

\phkoushiki{%
\phk{音は空気・水・固体の中を伝わる。真空の中では伝わらない}{中学理科}{ph:波-9:A4}
\phk{音の速さは約\qty{340}{m/s}}{中学理科}{ph:波-9:A5}
\phk{音の速さ $V=331.5+0.6t$（$t$ はセルシウス温度）}{高校物理基礎}{ph:波-9:A17}
}
%% 整合（U-5・K-8）：小学校理科の「わかること」の行（音が出るとき物は震えている）をけずった（本文①はそのまま）。

\begin{hatena}
\noindent\begin{minipage}[c]{0.56\linewidth}\raggedright
遠くの空で、ピカッと稲妻が光りました。それから「1、2、3」と数えたころに、ゴロゴロと雷の音が聞こえてきます。

光と音は、同じ雷から同時に出たはずです。なぜ、音だけが少しおくれて届くのでしょう。おくれた時間から、雷までの距離がわかるでしょうか。
\end{minipage}\hfill
\begin{minipage}[c]{0.40\linewidth}\centering
\begin{tikzpicture}[x=0.8mm, y=0.8mm]
  \fill[bkGray!30] (2,30) .. controls (0,36) and (8,40) .. (13,37) .. controls (16,43) and (27,42) .. (28,36) .. controls (34,37) and (36,30) .. (30,28) -- (4,28) .. controls (0,28) and (0,30) .. (2,30) -- cycle;
  \draw[bkGray!80!black, line width=0.6pt] (2,30) .. controls (0,36) and (8,40) .. (13,37) .. controls (16,43) and (27,42) .. (28,36) .. controls (34,37) and (36,30) .. (30,28) -- (4,28) .. controls (0,28) and (0,30) .. (2,30) -- cycle;
  \filldraw[fill=bkYellow, draw=bkOrange!80!black, line width=0.6pt] (16,28) -- (12,18) -- (16,18.5) -- (11,6) -- (21,20) -- (17,19.5) -- (21,28) -- cycle;
  \fill[bkGreen!25] (0,0) rectangle (60,3);
  \draw[bkInk, line width=0.6pt] (0,3) -- (60,3);
  \filldraw[fill=bkOrange!25, draw=bkInk, line width=0.6pt] (48,3) rectangle (58,11);
  \filldraw[fill=bkRed!40, draw=bkInk, line width=0.6pt] (46.5,11) -- (53,16.5) -- (59.5,11) -- cycle;
  \draw[rxj arr, bkOrange!85!black] (23,22) -- (44,13);
  \node[rxj lab, anchor=south west, text=bkOrange!70!black] at (30,20) {光（実線）：すぐ};
  \draw[rxj arr, bkBlue!80!black, dashed] (21,14) -- (44,7);
  \node[rxj lab, anchor=north, text=bkBlue!70!black] at (32,-1) {音（破線）：3 s 後};
\end{tikzpicture}
\end{minipage}
\end{hatena}

\noindent 熱-4 では、音を「空気の小さな圧力の波」として、その速さを導きました。この章では、その音を\textbf{はかる}ことを学びます。音の速さは、空気の中・水の中・鉄の中でどう違うのか。音の「大きさ」は何で決まり、どうやって数にするのか。うるさい音を小さくするには、どうすればよいのか。身の回りの騒音や、コンサートホールの響きまで、波の式で考えていきます。

%% ---------------- やってみよう・思い出そう ----------------
\begin{yattemiyou}[こだまで音の速さをはかろう]
\textbf{用意するもの}\par
校舎などの大きな平らな壁が見える広い場所、巻き尺（または歩数）、ストップウォッチ（スマホでよい）、手をたたく人と時間をはかる人
\begin{enumerate}[label={\bkcircled[bkGreen]{\arabic*}}, leftmargin=2.4\zw, itemsep=1pt, topsep=2pt]
\item 壁から\qty{50}{m}ほど離れて立ち、壁までの距離をはかる。周りに、ほかの建物や車がない所を選ぶ。
\item 手を1回大きくたたき、壁からはね返ってきた音（こだま）が聞こえるまでの時間をはかる。短くてはかりにくいときは、こだまが聞こえた瞬間に次をたたく、を10回続けて、10回分の時間をはかる。
\item その日の気温も調べておく。
\end{enumerate}
\end{yattemiyou}

\noindent 音は壁まで行って、もどってきます。だから、音が進んだ道のりは、壁までの距離の2倍です。2倍の距離を時間でわれば、音の速さが出ます（\bkreidairef{2}）。人が手をたたいてから時計を止めるまでには、0.1〜0.2秒のおくれがあるので、10回分を続けてはかると、ずれが小さくなります。

\begin{omoidasou}
\bkrecall{熱-4}{気体の中の音の速さ $a=\sqrt{\gamma R_{\mathrm s}T}$（⑥）。\qty{20}{\celsius}の空気で\qty{343}{m/s}、ヘリウムで\qty{1008}{m/s}。この巻では、波の速さを $c$ と書くので、音の速さも $c$ と書きます（高校の教科書の $V$ のこと）。}
\bkrecall{熱-3}{すばやく押す（断熱の）空気の体積弾性率 $K_{\mathrm s}=\gamma p$（⑧）。1気圧の空気で約\qty{0.14}{MPa}。}
\bkrecall{波-5}{波の速さ・振動数・波長の関係 $c=f\lambda$（②）。流体の中の音の速さ $c=\sqrt{K/\rho}$（⑤。小さな部分の運動方程式から、弦と同じ形の波動方程式を導いた）。}
\bkrecall{力-8}{体積弾性率 $K$：体積が縮む割合と、圧力の増え方の比。水は約\qty{2.2}{GPa}。}
\bkrecall{3-12}{デシベル＝$10\log_{10}$（強さの比）＝$20\log_{10}$（電圧の比）（⑦）。2倍は約3\,dB、10倍は10\,dB。}
\bkrecall{5-17}{$x$ が小さいとき $(1+x)^r\approx1+rx$（⑥の表）。$r=\tfrac12$ とすると $\sqrt{1+x}\approx1+\tfrac12x$。}
\end{omoidasou}

%% ---------------- 物理の道具 ----------------
\Needspace{16\baselineskip}
\bksection{doguu}{物理の道具}

\bksubsection{1}{音は、空気の圧力のゆれ}

\leavevmode\label{ph:波-9:sec1}たいこをたたくと、皮が震えます。指で軽くさわると、ビリビリと震えているのがわかります。強くたたくと、皮は大きく震え、音も大きくなります。小学校理科で学んだ「音が出るとき物は震えている。大きい音ほど、震え方が大きい」です。では、震える皮は、どうやって耳まで音を届けるのでしょうか。

スピーカーの膜（コーン）で考えます（図\ref{ph:波-9:fig-mitsu}）。膜が前へ出ると、すぐ前の空気を押し縮めます。押し縮められた空気は、さらに前の空気を押します。膜が後ろへ下がると、前の空気は引きのばされて薄くなります。こうして、空気が\textbf{こい所（密）}と\textbf{薄い所（疎）}が、交互に前へ進んでいきます。空気の粒そのものは、その場で前後に少し動くだけです。進んでいくのは「こい・薄い」という\textbf{形}です。空気が動く向きと、波が進む向きが同じなので、音は\textbf{縦波}です（波-5 ①）。

\begin{bkfigure}
\begin{tikzpicture}[x=1mm, y=1mm]
  % スピーカー
  \filldraw[fill=bkGray!40, draw=bkInk, line width=0.7pt] (0,12) rectangle (5,28);
  \filldraw[fill=bkGray!15, draw=bkInk, line width=0.7pt] (5,16) -- (11,6) -- (11,34) -- (5,24) -- cycle;
  \node[rxj lab, anchor=north] at (6,5) {スピーカー};
  \draw[rxj arr] (12.5,20) -- (17,20);
  \node[rxj lab, anchor=south] at (15,21) {押す};
  % 空気の粒
  \rxjair{20}{128}{20}{24}{3.6}{36}
  % 密・疎のしるし
  \foreach \x/\t in {20/密,38/疎,56/密,74/疎,92/密,110/疎,128/密} {\node[rxj lab, anchor=south] at (\x,33.5) {\t};}
  % 圧力のグラフ
  \draw[bkInk, line width=0.6pt, -{Stealth[length=3pt]}] (18,-12) -- (132,-12) node[rxj lab, anchor=west] {$x$};
  \draw[bkInk, line width=0.6pt, -{Stealth[length=3pt]}] (18,-20) -- (18,-3) node[rxj lab, anchor=south west] {音圧 $p$};
  \draw[bkBlue!80!black, line width=1.3pt, domain=20:128, samples=120] plot (\x, {-12+6*cos(10*(\x-20))});
  \node[rxj lab, anchor=east] at (17.2,-12) {0};
  \draw[rxj dim] (20,-21.5) -- (56,-21.5);
  \node[rxj lab, anchor=north] at (38,-22.5) {波長 $\lambda$};
  \draw[rxj arr, bkRed!75!black] (100,-23) -- (126,-23);
  \node[rxj lab, anchor=north] at (113,-24) {速さ $c$ で進む};
\end{tikzpicture}
\bkfigcaption{スピーカーの膜が空気を押したり引いたりすると、空気のこい所（密）と薄い所（疎）が交互にでき、前へ進む。下のグラフは、そのときの音圧（大気圧からのずれ）。密の所で正、疎の所で負。}
\label{ph:波-9:fig-mitsu}
\end{bkfigure}

空気がこい所では、圧力が大気圧より少し高く、薄い所では少し低くなります。この「大気圧からのずれ」を\textbf{音圧}といい、$p$ と書きます。\bkidx{おんあつ}{音圧}単位は圧力と同じ Pa です。音圧は、とても小さな量です。ふつうの会話の音圧は約\qty{0.02}{Pa}で、大気圧（約\qty{101000}{Pa}）の500万分の1くらいしかありません。耳は、このわずかな圧力のゆれを感じとっているのです（表\ref{ph:波-9:tab-onatsu}）。

\begin{bkcenter}
\begin{bktable}{colspec={X[2.2,l]Q[c,m]Q[c,m]}, width=0.92\linewidth}
音 & 音圧 $p$〔Pa〕 & 大気圧の何分の1 \\
聞こえ始める、いちばん小さな音 & 0.000\,02 & 約50億分の1 \\
ふつうの会話（\qty{1}{m}の所） & 0.02 & 約500万分の1 \\
電車が通るときのガード下 & 2 & 約5万分の1 \\
耳が痛くなるほどの音 & 20〜60 & 約5000分の1〜 \\
大気圧（比べるため） & \num{101325} & 1 \\
\end{bktable}
\bktabcaption{音圧の大きさ（実効値。⑥で説明する）。いちばん小さな音といちばん大きな音で、100万倍以上も違う。}
\label{ph:波-9:tab-onatsu}
\end{bkcenter}

\noindent\textbf{音圧は、膜が押す速さで決まる}\par\nobreak
熱-4 ⑥では、ピストンを速さ $\dd v$ で押したとき、押し縮められた所の圧力が $\dd p=\rho a\,\dd v$ だけ上がることを導きました。ピストンを、スピーカーの膜に置きかえます。膜が速さ $u$ で空気を押すと、膜のすぐ前の空気も速さ $u$ で動き、そこの音圧は
\pheq{p=\rho_0\,c\,u}{Pa ＝ kg/m$^3$ × m/s × m/s}
\label{ph:波-9:A1}になります（$\rho_0$ は空気の密度、$c$ は音の速さ）。膜が大きく震えるほど、膜の速さ $u$ は大きくなり、音圧 $p$ も大きくなります。これが「大きい音ほど、震え方が大きい」の式です。同じ式は、空気の中のどこでも成り立ちます。音が通っていくとき、空気の粒は速さ $u=p/(\rho_0c)$ で前後に動いているのです。

\begin{phyomi}[$p=\rho_0cu$]
「音圧は、空気の密度と音の速さと、空気が動く速さのかけ算」。\qty{20}{\celsius}の空気では $\rho_0c=1.204\times343\approx\qty{413}{Pa.s/m}$ なので、空気を\qty{1}{mm/s}で押すごとに、約\qty{0.41}{Pa}の音圧になります。$\rho_0c$ は、その物質が「どれだけ押し返してくるか」を表す数で、\textbf{音響インピーダンス}と呼ばれます（詳しくは波-7）。
\end{phyomi}
\begin{phmijika}[空気の粒は、ほとんど動かない]
\qty{1}{Pa}の音（⑥で計算すると94\,dB。かなり大きな音）でも、空気が動く速さは $1\div413\approx\qty{2.4}{mm/s}$ です。\qty{1000}{Hz}の音なら、空気は1秒に1000回往復するので、空気のずれ（実効値）は1万分の4 mm ほどです（\bkreidairef{3}）。聞こえ始めの\qty{20}{\micro Pa}では、動く幅は原子1個の大きさの数十分の1しかありません。耳は、それほどわずかな動きを聞き分けています。
\end{phmijika}

\noindent\textbf{聞こえる音の高さと、波長}\par\nobreak
1秒に震える回数（振動数 $f$）が多いほど、音は高く聞こえます。人の耳に聞こえるのは、およそ\qty{20}{Hz}から\qty{20000}{Hz}までです。これより高い音を\textbf{超音波}、低い音を\textbf{超低周波音}といいます。\bkidx{ちょうおんぱ}{超音波}波長は $\lambda=c/f$（波-5 ②）なので、空気の中では
\[ \lambda=\frac{343}{20}\approx\qty{17}{m}\quad\text{から}\quad\lambda=\frac{343}{\num{20000}}\approx\qty{17}{mm} \]
までです（\bkreidairef{15}）。低い音の波長は、建物の大きさくらい、高い音の波長は、指の太さくらいです。この差が、低い音は塀の向こうへ回りこみやすい（回折。波-8）、高い音はまっすぐ進みやすい、という違いを生みます。

%% ---------------------------------------------------------------
\bksubsection{2}{雷の音はなぜおくれる：光と音の速さ}

\leavevmode\label{ph:波-9:sec2}\bkhatenaref に答えましょう。稲妻の光と雷の音は、同じ瞬間に出ます。光は1秒に約30万 km（\qty{3.0e8}{m/s}）進みます。\qty{1}{km}先からなら、届くのにかかる時間は $1000\div(3.0\times10^8)\approx\qty{3.3e-6}{s}$、つまり100万分の3秒ほどで、私たちには「すぐ」です。一方、音は1秒に約\qty{340}{m}しか進みません。\qty{1}{km}進むのに約3秒かかります。

\textbf{はてな？の答え}：音の速さは、光の速さの約100万分の1しかないので、同じ所から同時に出ても、音のほうがおくれて届くのです。光ってから音が聞こえるまでの時間を $t$ とすると、光の時間は無視できるので、雷までの距離は
\pheq{(\text{距離})=c\,t\approx340\times t}{m ＝ m/s × s}
です。3秒なら約\qty{1.0}{km}です（\bkreidairef{1}）。「秒数を3でわると、およその km」と覚えておくと便利です（図\ref{ph:波-9:fig-kaminari}）。

\begin{bkfigure}
\begin{tikzpicture}
\begin{axis}[bkplot, width=100mm, height=46mm, xmin=0, xmax=10.5, ymin=0, ymax=3.8, xtick={1,2,3,4,5,6,7,8,9,10}, ytick={1,2,3}, xlabel={光ってからの時間 $t$（s）}, ylabel={雷までの距離（km）}, axis x line*=bottom, axis y line*=left, every axis x label/.style={at={(axis description cs:0.5,-0.18)}, anchor=north, font=\sffamily\small}]
  \addplot[bkcurve, domain=0:10.3] {0.340*x};
  \fill[bkOrange] (axis cs:3,1.02) circle (2pt); \fill[bkOrange] (axis cs:6,2.04) circle (2pt); \fill[bkOrange] (axis cs:9,3.06) circle (2pt);
  \node[rxj lab, anchor=west] at (axis cs:3.2,0.85) {3 s で約 1 km};
  \node[rxj lab, anchor=west] at (axis cs:6.2,1.85) {6 s で約 2 km};
\end{axis}
\end{tikzpicture}
\bkfigcaption{光ってから音が聞こえるまでの時間と、雷までの距離（音の速さ\qty{340}{m/s}）。グラフは原点を通る直線で、傾きが音の速さ。}
\label{ph:波-9:fig-kaminari}
\end{bkfigure}

花火や、運動会のピストル、遠くでくいを打つ工事でも、見えてから音が聞こえるまでに、少しおくれがあります。ピストルの音でスタートの時間をはかると、\qty{100}{m}離れた所では、約0.3秒おそくはかってしまいます。陸上競技の公式の記録で、ピストルの煙を見てストップウォッチを押していたのは、このためです。

\begin{phmodel}[雷までの距離]
\phignore{光が届くまでの時間（100万分の数秒）。風。気温による音の速さの違い。}
\phassume{音はまっすぐ、一定の速さ\qty{340}{m/s}で進む。稲妻の光った所と、音の出た所は同じ。}
\phlimit{稲妻は何 km も長いので、近い所と遠い所から出た音が、ずれて届く（ゴロゴロと長く続くわけ）。気温が高い・低い（⑤）。}
\end{phmodel}

%% ---------------------------------------------------------------
\bksubsection{3}{音の速さは、縮みにくさと密度で決まる}

\leavevmode\label{ph:波-9:sec3}波-5 ⑤では、管の中の水や空気の小さな部分に運動方程式を立てて、弦と同じ形の\textbf{波動方程式}を導き、流体の中の音の速さが
\[ c=\sqrt{\frac{K}{\rho_0}} \]
\label{ph:波-9:sqrtK}で決まることを確かめました（$K$ は体積弾性率、$\rho_0$ は密度。波-5 ⑤では $\rho$ と書いた）。押し返す力が強い（$K$ が大きい）ほど、となりへ早く伝わって速く、重い（$\rho_0$ が大きい）ほど、動き出すのに時間がかかって遅くなります（波-5 ⑤の「もどそうとする力の強さ÷重さ」の平方根）。この章では、この式を、空気・水・鉄の中の音と、気温による音の速さの変わり方に当てはめます。

\noindent\textbf{気体に当てはめる}\par\nobreak
音は1秒に何百回も押し縮めるので、熱が移る時間はなく、断熱で縮みます。気体の縮みにくさは、熱-3 ⑧の断熱の体積弾性率 $K_{\mathrm s}=\gamma p_0$ です（波-5 ⑤の $\sqrt{\gamma p/\rho}$）。理想気体では $p_0/\rho_0=R_{\mathrm s}T$ なので
\[ c=\sqrt{\frac{\gamma p_0}{\rho_0}}=\sqrt{\gamma R_{\mathrm s}T} \]
と、絶対温度 $T$ だけで書けます。熱-4 ⑥の $a=\sqrt{\gamma R_{\mathrm s}T}$ と同じ式です。\qty{20}{\celsius}の空気で約\qty{343}{m/s}です。温度一定で計算したニュートンの値（約\qty{290}{m/s}）が小さすぎた話は、熱-4 ⑥のコラムにあります。

\Needspace{4\baselineskip}
\noindent\textbf{真空の中では、音は伝わらない}\par\nobreak
\label{ph:波-9:A4}音の波動方程式は、媒質の小さな部分の運動方程式から生まれました（波-5 ⑤）。押し返す物も、動く物もない真空では、$K$ も $\rho_0$ もなく、式そのものが立ちません。だから、音は真空の中を伝わりません。中学理科の「音は空気・水・固体の中を伝わる。真空の中では伝わらない」は、このことです。容器の中に鳴っているブザーを入れ、ポンプで空気を抜いていくと、音はだんだん小さくなり、ほとんど聞こえなくなります。宇宙空間で、となりの宇宙船が爆発しても、音は聞こえません（光は届きます）。

\Needspace{10\baselineskip}
\begin{phflow}
\phflowrow{$c=\sqrt{K/\rho_0}$\\（波-5 ⑤）}{気体・断熱\\$K_{\mathrm s}=\gamma p_0$}{$c=\sqrt{\gamma R_{\mathrm s}T}$\\（熱-4 ⑥と同じ）}
\phflowrow{$c=\sqrt{\gamma R_{\mathrm s}(273.15+\theta)}$}{$\theta$ が小さい\\1次近似（5-17⑥）}{$c\approx331.5+0.6\,\theta$\\（⑤）}
\end{phflow}

%% ---------------------------------------------------------------
\bksubsection{4}{水の中・鉄の中の音の速さ}

\leavevmode\label{ph:波-9:sec4}$c=\sqrt{K/\rho_0}$ は、液体にもそのまま使えます。固体では、縮みにくさ $K$ を、形に合った弾性の定数に置きかえた形になります（波-5 ⑤）。

\noindent\textbf{液体}\par\nobreak
水の体積弾性率は約\qty{2.2}{GPa}（力-8）で、空気の断熱の値（\qty{0.14}{MPa}）の約1万5000倍も縮みにくいのです。密度は約\qty{998}{kg/m^3}で、空気の約800倍重い。縮みにくさのほうがずっと大きく効くので
\[ c=\sqrt{\frac{2.2\times10^9}{998}}\approx\qty{1480}{m/s} \]
と、空気の4倍以上の速さになります（\bkreidairef{4}）。海水は塩がとけている分だけ縮みにくく、約\qty{1500}{m/s}です。船の魚群探知機や、海の深さをはかるソナーは、この速さを使っています（波-10）。

\noindent\textbf{固体}\par\nobreak
固体は、押し縮めるだけでなく、ずらす（せん断の）力にも逆らいます。そのため、固体の中の音の速さは、形によって変わります。
\begin{itemize}[leftmargin=1.2\zw, itemsep=2pt, topsep=2pt, label={\color{bkTheme}●}]
\item \textbf{細い棒}（レールや針金）：棒は、押されると横に少しふくらめるので、縮みにくさはヤング率 $E$ です。$c=\sqrt{E/\rho_0}$（波-5 ⑤）。鋼（$E=\qty{205}{GPa}$、$\rho_0=\qty{7850}{kg/m^3}$）で約\qty{5100}{m/s}（力-16 の値）。
\item \textbf{広い固体の中}（厚い鋼の板の中、地面の中）：横にふくらめないので、縮みにくさは $E$ より大きくなり、鋼で約\qty{5900}{m/s}です（もっと知りたい）。
\item \textbf{横波}：固体の中では、ずらす向きのゆれ（横波）も伝わります。速さは $\sqrt{G/\rho_0}$（$G$ は横弾性係数。力-8 ④）で、鋼で約\qty{3200}{m/s}。液体と気体は、ずらしても元にもどる力がないので、横波は伝わりません。
\end{itemize}

\begin{bkfigure}
\begin{tikzpicture}[x=1mm, y=1mm]
  \foreach \v in {1000,2000,...,6000} {\draw[bkLine, line width=0.3pt] ({\v*0.015},0) -- ({\v*0.015},51);}
  \foreach \v in {0,1000,2000,...,6000} {\node[rxj lab, anchor=north] at ({\v*0.015},-1.5) {\v};}
  \foreach \i/\v/\t in {1/331/空気（0 ℃）,2/343/空気（20 ℃）,3/1008/ヘリウム（20 ℃）,4/1480/水（20 ℃）,5/1500/海水,6/5100/鋼（細い棒）,7/5900/鋼（広い固体）} {%
    \filldraw[fill=bkBlue!45, draw=bkInk, line width=0.5pt] (0,{\i*7-5.2}) rectangle ({\v*0.015},{\i*7-0.8});
    \node[rxj lab, anchor=east] at (-1.5,{\i*7-3}) {\t};
    \node[rxj lab, anchor=west] at ({\v*0.015+1.2},{\i*7-3}) {\v};}
  \draw[bkInk, line width=0.8pt, -{Stealth[length=4pt]}] (0,0) -- (98,0);
  \draw[bkInk, line width=0.8pt] (0,0) -- (0,51);
  \node[rxj lab, font=\sffamily\small, anchor=north] at (45,-6.5) {音の速さ（m/s）};
\end{tikzpicture}
\bkfigcaption{いろいろな物の中の音の速さ（縦波）。固体は、気体の十数倍も速い。}
\label{ph:波-9:fig-hayasa}
\end{bkfigure}

\begin{phmijika}[線路に耳をあてると、汽車の音が早く聞こえる]
昔の子どもの遊びに、線路に耳をあてて、遠くの汽車の近づく音を聞く、というものがありました（とても危ないので、まねしてはいけません）。鋼のレールの中では、音が空気の15倍ほど速く進みます。\qty{1}{km}先の音は、レールを通ると約0.2秒、空気を通ると約2.9秒で届くので（\bkreidairef{5}）、レールの音のほうが先に聞こえるのです。レールの中では音があまり弱まらないことも、遠くまで聞こえるわけです。
\end{phmijika}

\begin{phbaisuru}{縮みにくさ $K$}
音の速さは $\sqrt2\approx1.41$ 倍になります。密度を2倍にすると、$1/\sqrt2\approx0.71$ 倍に遅くなります。ヘリウムが空気より約3倍速いのは、縮みにくさ（$\gamma p_0$）は空気の約1.2倍とあまり変わらないのに、密度が約7分の1しかないからです。ヘリウムを吸って声を出すと、のどや口の中の空気の柱の固有振動数（波-6 ④）が上がり、声の響きが高く変わります。
\end{phbaisuru}

\begin{motto}[4]
\phboxhead{広い固体の中の縦波と横波：地震の P 波と S 波}
広い固体の中で縦波が進むとき、横には少しもふくらめません。そのため、縮みにくさは体積弾性率 $K$ に、ずらす力への強さの分 $\tfrac43G$ が加わった $K+\tfrac43G$ になります。鋼（$E=\qty{205}{GPa}$、ポアソン比0.3）では $K\approx\qty{171}{GPa}$、$G\approx\qty{79}{GPa}$ で
\[ c_{\text{縦}}=\sqrt{\frac{K+\tfrac43G}{\rho_0}}=\sqrt{\frac{(171+105)\times10^9}{7850}}\approx\qty{5930}{m/s} \]
\[ c_{\text{横}}=\sqrt{\frac{G}{\rho_0}}=\sqrt{\frac{79\times10^9}{7850}}\approx\qty{3170}{m/s} \]
です。地震のときも、地面の中を縦波（P 波）と横波（S 波）が伝わります。縦波のほうが速いので、先に小さなゆれ（初期微動）が来て、あとから大きなゆれが来ます。その時間の差から、震源までの距離がわかります（波-5 ①）。
\end{motto}

%% 修3-4：章末のページの空き（コラムが入りきらず下の空きが4割）を直すため、コラムを水の中の音の速さの節（④）の終わりへ移した
\begin{column}{音の速さをはかった人たち}
空気の中の音の速さは、17世紀から18世紀にかけて、遠くで撃った大砲の光が見えてから音が聞こえるまでの時間をはかって、求められ（\bkreidairef{1}と同じ考え方）、およそ\qty{330}{m/s}という値が得られました。

水の中の音の速さは、1826年、スイスのコラドンとフランスのシュトゥルムが、ジュネーブ湖ではかりました。一方の船で水の中の鐘を鳴らすと同時に火薬を光らせ、\qty{10}{km}あまり離れたもう一方の船で、水の中に入れたラッパのような筒に耳をあてて、光が見えてから鐘の音が届くまでの時間をはかったのです。答えは約\qty{1435}{m/s}で、今の精密な測定による約\qty{8}{\celsius}の水の値（約\qty{1440}{m/s}）とよく合っています。水は温度が下がると少し縮みやすくなるので、\qty{20}{\celsius}の水（約\qty{1480}{m/s}）より遅いのです。
\end{column}

%% ---------------------------------------------------------------
\bksubsection{5}{温度で変わる音の速さ}

\leavevmode\label{ph:波-9:sec5}気体の音の速さ $c=\sqrt{\gamma R_{\mathrm s}T}$ は、絶対温度 $T$ の平方根に比例します。高校物理基礎では、空気の音の速さを
\[ V=331.5+0.6t\qquad(t\text{ はセルシウス温度〔℃〕}) \]
と習います。平方根の式から、この1次式（直線の式）を導きましょう。この本では、セルシウス温度を $\theta$ と書きます（熱-1）。

\begin{enumerate}[label={\bkcircled[bkBlue]{\arabic*}}, leftmargin=2.4\zw, itemsep=2pt, topsep=2pt]
\item $T=273.15+\theta$ を入れて、\qty{0}{\celsius}の音の速さ $c_0=\sqrt{\gamma R_{\mathrm s}\times273.15}$ をくくり出す。
\[ c=\sqrt{\gamma R_{\mathrm s}(273.15+\theta)}=c_0\sqrt{1+\frac{\theta}{273.15}} \]
\item $\theta$ が 273.15 よりずっと小さいとき、$\sqrt{1+x}\approx1+\tfrac12x$（5-17⑥）を使う。
\[ c\approx c_0\Bigl(1+\frac{\theta}{2\times273.15}\Bigr)=c_0+\frac{c_0}{546.3}\,\theta \]
\item 空気（$\gamma=1.4$、$R_{\mathrm s}=\qty{287}{J/(kg.K)}$）では、$c_0=\sqrt{1.4\times287\times273.15}\approx\qty{331.3}{m/s}$、傾きは $331.3\div546.3\approx0.606$。
\end{enumerate}
\pheq{c\approx331.3+0.606\,\theta\ \approx\ 331.5+0.6\,\theta}{m/s ＝ m/s ＋ m/(s·K) × K}
\label{ph:波-9:A17}高校の式が出てきました。0.6 は「\qty{0}{\celsius}の音の速さを、$2\times273.15$ でわった数」だったのです。高校の 331.5 は、乾いた空気ではかった\qty{0}{\celsius}の値です（$\gamma$ と $R_{\mathrm s}$ を3けたで使った式の値 331.3 と、0.2 しか違いません）。

\begin{bkfigure}
\begin{tikzpicture}
\begin{axis}[bkplot, width=108mm, height=52mm, xmin=-45, xmax=45, ymin=300, ymax=362, xtick={-40,-20,0,20,40}, ytick={306,320,331,343,355}, xlabel={$\theta$（℃）}, ylabel={$c$（m/s）}, axis x line*=bottom, axis y line*=left]
  \addplot[bkcurve, domain=-42:42] {sqrt(1.4*287*(x+273.15))};
  \addplot[bkOrange, line width=1.2pt, dashed, domain=-42:42] {331.5+0.6*x};
  \node[rxj lab, anchor=west, text=bkBlue!80!black] at (axis cs:2,318) {実線：$\sqrt{\gamma R_{\mathrm s}T}$};
  \node[rxj lab, anchor=west, text=bkOrange!80!black] at (axis cs:2,309) {破線：$331.5+0.6\theta$};
  \fill[bkInk] (axis cs:-40,306.1) circle (1.8pt); \fill[bkInk] (axis cs:0,331.3) circle (1.8pt); \fill[bkInk] (axis cs:20,343.2) circle (1.8pt); \fill[bkInk] (axis cs:40,354.7) circle (1.8pt);
\end{axis}
\end{tikzpicture}
\bkfigcaption{空気の音の速さと温度。実線は $\sqrt{\gamma R_{\mathrm s}T}$、破線は高校の1次式。$-40$〜\qty{40}{\celsius}で、差は\qty{1.5}{m/s}より小さい（例えば\qty{40}{\celsius}で354.7と355.5）。}
\label{ph:波-9:fig-ondo}
\end{bkfigure}

\begin{phyomi}[$c\approx331.5+0.6\theta$]
「気温が\qty{1}{\celsius}（\qty{1}{K}）上がるごとに、音は\qty{0.6}{m/s}ずつ速くなる」。\qty{0}{\celsius}で331.5、\qty{15}{\celsius}で340.5、\qty{20}{\celsius}で343.5、\qty{30}{\celsius}で349.5です。\label{ph:波-9:A5}中学理科の「音の速さは約\qty{340}{m/s}」は、\qty{15}{\celsius}くらいの値です。光の速さ（約\qty{3.0e8}{m/s}、②）の約100万分の1です。
\end{phyomi}
\begin{phbaisuru}{気温の差（$\theta$ の変わり方）}
夏（\qty{30}{\celsius}）と冬（\qty{0}{\celsius}）の差 \qty{30}{K}で、音の速さの差は $0.6\times30=\qty{18}{m/s}$。差を2倍の\qty{60}{K}にすると、速さの差も2倍の\qty{36}{m/s}です。1次式なので、温度の差に比例します。ただし、$\theta$ が\qty{100}{\celsius}をこえるような大きな範囲では、直線と平方根の曲線が離れていくので、もとの式を使います。
\end{phbaisuru}

\begin{phmijika}[夜は、遠くの音がよく聞こえる]
晴れた夜は、地面が冷えて、地面に近い空気ほど冷たく、上の空気ほど暖かくなることがあります。上のほうが音が速いので、上へ向かった音は、速い所で前へ進みすぎて、下へ曲がってきます（屈折。波-8）。昼に聞こえない遠くの電車の音や、踏切の音が、夜によく聞こえるのは、このためです。
\end{phmijika}

\begin{motto}[3]
\phboxhead{音の速さは、分子の速さと同じくらい}
熱-14 ⑤で、\qty{20}{\celsius}の窒素の分子が、平均で約\qty{471}{m/s}で飛びまわっていることを学びました。音の速さ\qty{343}{m/s}は、その約0.73倍です。音は、分子がとなりの分子にぶつかって押していくことで伝わるので、分子の速さより速くは伝わりません。温度が上がると分子が速くなり、音も速くなります。音の速さが $\sqrt T$ に比例するのは、分子の速さが $\sqrt T$ に比例するのと同じ理由なのです。気圧が変わっても、分子の速さは変わらないので、音の速さも変わりません。
\end{motto}

%% ---------------------------------------------------------------
\bksubsection{6}{音の大きさを数にする：デシベル}

\leavevmode\label{ph:波-9:sec6}音圧は、たえず正と負に変わっているので、平均すると0になってしまいます。そこで、音の大きさには、2乗して平均し、平方根をとった値（\textbf{実効値}）を使います。正弦波の形の音なら、実効値は、音圧のいちばん大きな値（振幅）の $1/\sqrt2\approx0.71$ 倍です。家の電気の「100 V」が実効値なのと同じです（3-10）。この章の音圧は、ことわらないかぎり実効値です。

\noindent\textbf{音の強さ}\par\nobreak
音は、エネルギーを運びます。音の進む向きに垂直な\qty{1}{m^2}を、1秒に通るエネルギーを\textbf{音の強さ}といい、$I$ と書きます。\bkidx{おとのつよさ}{音の強さ}単位は W/m$^2$ です。音が通る面では、音圧 $p$ が、速さ $u$ で動く空気を押しています。仕事率は「力×速さ」なので、\qty{1}{m^2}あたりでは $p\times u$ です。①の $u=p/(\rho_0c)$ を入れると
\pheq{I=p\,u=\frac{p^2}{\rho_0c}}{W/m$^2$ ＝ Pa$^2$÷(Pa·s/m)}
です（2乗の平均なので、$p$ は実効値。詳しくは波-7）。強さは、音圧の\textbf{2乗}に比例します。

\noindent\textbf{レベル：けたで数える}\par\nobreak
耳が聞きとれる音圧は、表\ref{ph:波-9:tab-onatsu}のように、100万倍以上も違います。強さにすると、1兆倍以上です。そこで、3-12⑦で学んだ\textbf{デシベル}（dB）で、基準の何倍かを「けた」で数えます。音の強さの基準を、聞こえ始めの強さに近い $I_0=\qty{1e-12}{W/m^2}$ として
\pheq{L_I=10\log_{10}\frac{I}{I_0}}{dB}
を\textbf{音の強さのレベル}といいます。\bkidx{おとのつよさのれべる}{音の強さのレベル}強さは音圧の2乗に比例するので、音圧で書くと、2乗が対数の前の2になり
\pheq{L_p=10\log_{10}\frac{p^2}{p_0^{\,2}}=20\log_{10}\frac{p}{p_0},\qquad p_0=\qty{20}{\micro Pa}}{dB}
\label{ph:波-9:Lp}となります。これを\textbf{音圧レベル}といいます。\bkidx{おんあつれべる}{音圧レベル}基準の $p_0=\qty{20}{\micro Pa}$ は、聞こえ始めの音圧です。$p_0^{\,2}/(\rho_0c)=(20\times10^{-6})^2\div413\approx\qty{0.97e-12}{W/m^2}$ で、$I_0$ とほぼ同じなので、空気の中では $L_p$ と $L_I$ は、ほぼ同じ数になります（差は0.2\,dB より小さい）。騒音計がはかるのは、音圧レベルです。

\begin{phyomi}[$L_p=20\log_{10}(p/p_0)$]
「音圧が10倍になるごとに20\,dB、2倍になるごとに約6\,dB 大きくなる」。強さで言えば「10倍で10\,dB、2倍で約3\,dB」です。音圧は強さの平方根なので、音圧の比には20をかけます（3-12⑦の電圧と同じ）。\qty{20}{\micro Pa}は0\,dB、\qty{0.02}{Pa}（会話）は60\,dB、\qty{1}{Pa}は94\,dB、\qty{20}{Pa}は120\,dB です（\bkreidairef{7}）。
\end{phyomi}

\begin{bkfigure}
\begin{tikzpicture}[x=1mm, y=1mm]
  \draw[bkInk, line width=1pt] (40,0) -- (40,72);
  \foreach \d in {0,20,...,120} {\draw[bkInk, line width=0.7pt] (38.5,\d*0.6) -- (41.5,\d*0.6); \node[rxj lab, anchor=west] at (43,\d*0.6) {\d\,dB};}
  \node[rxj lab, anchor=south, font=\sffamily\small\bfseries] at (50,74) {音圧レベル};
  \node[rxj lab, anchor=south, font=\sffamily\small\bfseries] at (14,74) {音圧};
  \foreach \d/\p in {0/0.000\,02 Pa,40/0.002 Pa,60/0.02 Pa,94/1 Pa,120/20 Pa} {\node[rxj lab, anchor=east] at (36,\d*0.6) {\p}; \draw[bkGray, line width=0.4pt] (36.3,\d*0.6) -- (38.3,\d*0.6);}
  \foreach \d/\t in {0/聞こえ始める音,30/ささやき声,40/静かな図書館,60/ふつうの会話,90/工場の機械の近く,100/電車が通るガード下,120/耳が痛くなるほどの音} {\fill[bkOrange!85!black] (40,\d*0.6) circle (0.9); \draw[bkGray, line width=0.4pt] (57,\d*0.6) -- (62,\d*0.6); \node[rxj lab, anchor=west] at (63,\d*0.6) {\t};}
\end{tikzpicture}
\bkfigcaption{身近な音の、音圧レベルのものさし（3-12 の値）。20\,dB 上がるごとに、音圧は10倍になる。}
\label{ph:波-9:fig-monosashi}
\end{bkfigure}

\noindent\textbf{なぜ、けたで数えるのか}\par\nobreak
3-12⑦で見たように、耳は、強さの「けた」に近い感じ方をします。強さが10倍（10\,dB）になっても、10倍うるさくは聞こえず、「2倍くらい大きくなった」と感じる人が多いことが、たくさんの人に聞いてもらう実験で知られています。dB は、耳の感じ方にも合った数え方なのです。

\begin{motto}[3]
\phboxhead{耳が感じる「音の大きさ」：フォンとソーン}
騒音計の dB は、空気の圧力のゆれの大きさです。耳が感じる「大きさ」は、それとは少し違います。そこで、\qty{1000}{Hz}の音で何 dB の音と同じ大きさに聞こえるかを「フォン」で表します（\qty{1000}{Hz}の音なら、dB とフォンは同じ数）。⑧の等ラウドネス曲線は、同じフォンの点を結んだ線です。

さらに、40フォンの音の大きさを「1ソーン」とし、「2倍の大きさに聞こえる」音を2ソーン、と数えるやり方もあります。たくさんの人に聞いてもらう実験から、音が10フォン大きくなるごとに、ソーンはほぼ2倍になることがわかっています。50フォンは2ソーン、60フォンは4ソーン、70フォンは8ソーンです。強さが10倍になっても、大きさは2倍にしか感じないのです。
\end{motto}

%% ---------------------------------------------------------------
\bksubsection{7}{dB の足し算と、離れると小さくなる音}

\noindent\label{ph:波-9:sec7}\textbf{音を合わせる}\par\nobreak
同じ80\,dB の機械を2台同時に動かすと、何 dB になるでしょうか。$80+80=160$\,dB ではありません。2台の音は、互いに関係なくゆれているので（波の山と谷がそろわない）、合わせると\textbf{強さ}がたし算になります（3-12 のまちがえやすいところ）。手順は次のとおりです。
\begin{enumerate}[label={\bkcircled[bkBlue]{\arabic*}}, leftmargin=2.4\zw, itemsep=2pt, topsep=2pt]
\item それぞれの dB を、基準の何倍の強さかにもどす：$I/I_0=10^{L/10}$。
\item 強さをたす。
\item たした強さを、もう一度 dB にする：$L=10\log_{10}(\text{たした強さの比})$。
\end{enumerate}
\pheq{L=10\log_{10}\bigl(10^{L_1/10}+10^{L_2/10}+\cdots\bigr)}{dB}
80\,dB が2台なら、強さの比は $10^8+10^8=2\times10^8$ なので
\[ L=10\log_{10}(2\times10^8)=80+10\log_{10}2\approx\qty{83}{dB} \]
です。同じ音を2台にしても、3\,dB しか増えません。80\,dB と70\,dB なら
\[ L=10\log_{10}(10^8+10^7)=80+10\log_{10}1.1\approx\qty{80.4}{dB} \]
で、小さいほうの音は、ほとんど効きません（\bkreidairef{9}）。

\begin{bkcenter}
\begin{bktable}{colspec={X[2,l]Q[c]Q[c]Q[c]Q[c]Q[c]}, width=0.92\linewidth}
2つの音のレベルの差〔dB〕 & 0 & 3 & 6 & 10 & 15 \\
大きいほうに足す量〔dB〕 & 3.0 & 1.8 & 1.0 & 0.4 & 0.1 \\
\end{bktable}
\bktabcaption{2つの音を合わせたときの増え方。差が10\,dB 以上なら、小さいほうはほぼ無視できる。}
\label{ph:波-9:tab-tashi}
\end{bkcenter}

反対に、工場で「機械を動かすと85\,dB、止めると周りの音（背景の音）だけで82\,dB」のときは、強さを\textbf{引き算}して、機械だけの音を求めます。$10\log_{10}(10^{8.5}-10^{8.2})\approx\qty{82}{dB}$ です。背景の音との差が10\,dB 以上あれば、表\ref{ph:波-9:tab-tashi}のように背景の音が効くのは0.4\,dB 以下なので、ふつうはこのひき算をしません。

\noindent\textbf{離れると小さくなる：点の音源}\par\nobreak
周りに何もない広い所で、小さなスピーカーが四方八方へ音を出しているとします。波-7 ②で見たように、音源が1秒に出すエネルギーは、距離 $R$ の所では半径 $R$ の球面（面積 $4\pi R^2$）に広がるので、強さは距離の2乗に反比例します（\textbf{逆2乗の法則}）。音源が1秒に出す音のエネルギーを\textbf{音響パワー}といい、音響の本にならって、この章では $W$〔W〕と書きます（波-7 ②の $P$ と同じ量。単位のワット W とは別の文字）。波-7 ②の式は
\[ I=\frac{W}{4\pi R^2} \]
です。音響パワーも基準 $W_0=\qty{1e-12}{W}$ との比を dB にした\textbf{音響パワーレベル} $L_W=10\log_{10}(W/W_0)$ で表すと\bkidx{おんきょうぱわーれべる}{音響パワーレベル}
\[ L_I=10\log_{10}\frac{W}{4\pi R^2\,I_0}=L_W-10\log_{10}(4\pi)-10\log_{10}R^2 \]
$10\log_{10}(4\pi)\approx11$ なので（$R$ は m で入れる）
\pheq{L_p\approx L_W-20\log_{10}R-11}{dB}
です（\bkreidairef{10}）。地面の上に置いた機械なら、音は半分の球（半球）にだけ広がるので、$-11$ が $-8$ になります。

\noindent\textbf{線の音源}\par\nobreak
車がたくさん通る道路は、まっすぐな線から音が出ているように見えます。線から出た音は、半径 $R$ の円柱の面に広がります。円柱の面積は $R$ に比例するので、強さは $R$ に反比例します（波-7 ②）。

\begin{phbaisuru}{音源からの距離 $R$}
点の音源なら、強さは $1/4$ で、$10\log_{10}4=20\log_{10}2\approx6$ なので、6\,dB 下がります。線の音源なら、強さは $1/2$ で、3\,dB しか下がりません（\bkreidairef{11}）。\qty{10}{m}で70\,dB の音が、\qty{40}{m}では、点の音源で58\,dB、線の音源で64\,dB です。道路の騒音が、離れてもなかなか小さくならないのは、このためです。
\par\vspace{2pt}\centering
\begin{tikzpicture}
\begin{axis}[bkplot, xmode=log, width=104mm, height=50mm, xmin=8, xmax=330, ymin=36, ymax=77, xtick={10,20,40,80,160,320}, xticklabels={10,20,40,80,160,320}, ytick={40,46,52,58,64,70}, xlabel={距離 $R$（m）},ylabel={$L_p$（dB）}, axis x line*=bottom, axis y line*=left, every axis x label/.style={at={(axis description cs:0.5,-0.17)}, anchor=north, font=\sffamily\small}]
  \addplot[bkcurve, domain=10:320, samples=40] {70-20*log10(x/10)};
  \addplot[bkOrange, line width=1.5pt, dashed, domain=10:320, samples=40] {70-10*log10(x/10)};
  \node[rxj lab, anchor=south west, text=bkOrange!80!black] at (axis cs:20,71) {破線：線の音源（2倍で $-3$\,dB）};
  \node[rxj lab, anchor=north west, text=bkBlue!80!black] at (axis cs:10.5,51) {実線：点の音源（2倍で $-6$\,dB）};
\end{axis}
\end{tikzpicture}
\end{phbaisuru}

%% ---------------------------------------------------------------
\bksubsection{8}{耳の感じ方と A 特性}

\leavevmode\label{ph:波-9:sec8}同じ60\,dB の音でも、\qty{100}{Hz}の低い音は、\qty{1000}{Hz}の音より小さく聞こえます。人の耳は、およそ\qty{2000}{Hz}〜\qty{5000}{Hz}の音にいちばん敏感で、低い音や、とても高い音には鈍いのです。いろいろな高さの音を、\qty{1000}{Hz}の音と「同じ大きさに聞こえる」ように調べて結んだ線を、\textbf{等ラウドネス曲線}といいます。\bkidx{とうらうどねすきょくせん}{等ラウドネス曲線}

騒音計は、この耳の感じ方に近づけるため、はかった音を、高さごとに決まった量だけ弱めてから合わせます。その弱め方を\textbf{A 特性}といい、A 特性でまとめた音圧レベルを「dB(A)」（エー特性音圧レベル）と書きます。\bkidx{えーとくせい}{A 特性}A 特性の補正の量（dB）は、規格の式で決まっていて、\qty{1000}{Hz}で0、\qty{100}{Hz}で約 $-19$\,dB、\qty{4000}{Hz}で約 $+1$\,dB です（図\ref{ph:波-9:fig-atoku}）。

\begin{bkfigure}
\begin{tikzpicture}
\begin{axis}[bkplot, xmode=log, width=110mm, height=54mm, xmin=20, xmax=20000, ymin=-50, ymax=8, xtick={31.5,63,125,250,500,1000,2000,4000,8000,16000}, xticklabels={31.5,63,125,250,500,1k,2k,4k,8k,16k}, ytick={-40,-30,-20,-10,0}, xlabel={振動数（Hz）}, ylabel={補正（dB）}, axis x line*=bottom, axis y line*=left, every axis x label/.style={at={(axis description cs:0.5,-0.15)}, anchor=north, font=\sffamily\small}]
  \addplot[bkcurve, domain=21:20000, samples=120] {20*log10((12194^2*x^4)/((x^2+20.6^2)*sqrt((x^2+107.7^2)*(x^2+737.9^2))*(x^2+12194^2)))+2.00};
  \fill[bkOrange] (axis cs:100,-19.1) circle (2pt); \fill[bkOrange] (axis cs:1000,0) circle (2pt); \fill[bkOrange] (axis cs:4000,1.0) circle (2pt);
  \node[rxj lab, anchor=west] at (axis cs:110,-22) {100 Hz：$-19.1$};
  \node[rxj lab, anchor=north west] at (axis cs:1000,-4) {1 kHz：0};
  \node[rxj lab, anchor=north] at (axis cs:5600,-12) {4 kHz：$+1.0$};
\end{axis}
\end{tikzpicture}
\bkfigcaption{A 特性の補正（規格の式で計算）。低い音ほど大きく弱める。耳が鈍い所を、騒音計も鈍くしている。}
\label{ph:波-9:fig-atoku}
\end{bkfigure}

たとえば、\qty{100}{Hz}で80\,dB の低い音は、A 特性で $80-19.1\approx61$\,dB(A) になります（\bkreidairef{12}）。大型トラックや空調の室外機の低い音は、dB(A) で見ると小さくなりますが、窓や壁を通りぬけやすく、体に響く感じがあるので、別に低い周波数の音として調べることもあります。

\begin{phmada}[耳の中には、音を大きくする「増幅器」がある]
耳の奥のうずまき管（かぎゅう）の中では、音の高さごとに、違う場所の膜が震えます。その膜の上には、\textbf{外有毛細胞}という細胞が並んでいて、音に合わせて自分で伸び縮みし、膜の震えを強めていると考えられています。とても小さな音が聞こえ、近い高さの音を聞き分けられるのは、このはたらきのおかげです。健康な耳からは、この細胞のはたらきで、とても小さな音が出てくることもわかっていて（耳音響放射）、生まれたばかりの赤ちゃんの聞こえの検査に使われています。

しかし、細胞の伸び縮みが、どのように膜の震えと組み合わさって、あれほど鋭い聞き分けになるのか、その全体のしくみはまだ解明されていません。音がないのに音が聞こえる「耳鳴り」がなぜ起きるのかも、耳と脳のどこで何が起きているのか、研究が続いています。大きな音を長く聞くと、外有毛細胞が傷つき、元にもどらないことがあります（騒音性難聴）。
\end{phmada}

%% ---------------------------------------------------------------
\bksubsection{9}{音を止める：遮音と質量則}

\leavevmode\label{ph:波-9:sec9}外の道路の音を、窓や壁でさえぎることを\textbf{遮音}といいます。\bkidx{しゃおん}{遮音}壁に当たった音の強さのうち、壁を通りぬけた強さの割合を $\tau$（タウ）とし、
\pheq{\mathrm{TL}=10\log_{10}\frac{1}{\tau}}{dB}
を\textbf{透過損失}といいます。\bkidx{とうかそんしつ}{透過損失}透過損失が30\,dB の窓なら、通りぬける音の強さは1000分の1で、外が75\,dB なら、部屋の中は約45\,dB です。

\noindent\textbf{重い壁ほど、音を通さない}\par\nobreak
音が壁に当たると、壁の音圧で壁全体が前後にゆすられます。ゆすられた壁は、反対側の空気を押して、そこに音を出します。壁が重いほど、同じ力でゆすられても動きにくく、反対側に出す音は小さくなります。壁の\qty{1}{m^2}あたりの質量（面密度）を $m$〔kg/m$^2$〕とします。壁を「ばねも、ねばりもない、重い板」とみて、壁の運動方程式と、両側の音圧の式を組み合わせると、垂直に当たる音の透過損失は
\pheq{\mathrm{TL}\approx20\log_{10}\frac{\pi fm}{\rho_0c}}{dB}
\label{ph:波-9:shitsuryo}になります（$\pi fm$ が $\rho_0c$ より十分大きいとき。log の中の $\pi fm/(\rho_0c)$ は単位のない数）。これを\textbf{質量則}といいます。\bkidx{しつりょうそく}{質量則}壁の重さ $m$ と、音の振動数 $f$ の積で決まります。

\begin{phyomi}[質量則]
「壁の重さが2倍になるか、音の振動数が2倍になると、透過損失は6\,dB 大きくなる」。$20\log_{10}2\approx6$ だからです。厚さ\qty{5}{mm}のガラス（$m\approx\qty{12.5}{kg/m^2}$）は、\qty{500}{Hz}で約34\,dB です（\bkreidairef{13}）。低い音ほど、通りぬけやすいことがわかります。
\end{phyomi}

\begin{bkfigure}
\begin{tikzpicture}
\begin{axis}[bkplot, xmode=log, width=106mm, height=54mm, xmin=100, xmax=4000, ymin=10, ymax=84, xtick={125,250,500,1000,2000,4000}, xticklabels={125,250,500,1k,2k,4k}, ytick={20,30,40,50,60,70}, xlabel={振動数 $f$（Hz）}, ylabel={TL（dB）}, axis x line*=bottom, axis y line*=left, every axis x label/.style={at={(axis description cs:0.5,-0.15)}, anchor=north, font=\sffamily\small}]
  \addplot[bkcurve, domain=100:4000, samples=40] {20*log10(pi*x*12.5/413)};
  \addplot[bkOrange, line width=1.5pt, dashed, domain=100:4000, samples=40] {20*log10(pi*x*25/413)};
  \addplot[bkGreen!70!black, line width=1.5pt, densely dotted, domain=100:1600, samples=40] {20*log10(pi*x*276/413)};
  \node[rxj lab, anchor=west, text=bkBlue!80!black] at (axis cs:1200,32) {実線：ガラス 5 mm（12.5 kg/m$^2$）};
  \node[rxj lab, anchor=west, text=bkOrange!80!black] at (axis cs:130,18) {破線：ガラス 10 mm（25 kg/m$^2$）};
  \node[rxj lab, anchor=west, text=bkGreen!50!black] at (axis cs:115,78) {点線：コンクリート 12 cm（276 kg/m$^2$）};
\end{axis}
\end{tikzpicture}
\bkfigcaption{質量則（垂直に当たる音）。振動数が2倍になるごと、重さが2倍になるごとに、6\,dB ずつ大きくなる。}
\label{ph:波-9:fig-shitsuryo}
\end{bkfigure}

\noindent\textbf{せいかくに言うと}\par\nobreak
本当の窓や壁では、音はいろいろな向きから当たるので、透過損失は、垂直に当たる音の質量則より5\,dB ほど小さくなるのがふつうです。また、ガラスのような硬い板には、板そのものが曲がって波打つ「板の波」と、空気の音がちょうど足並みをそろえて、音を通しやすくなる振動数（コインシデンス）があります。厚さ\qty{5}{mm}のガラスでは、数千 Hz のあたりで、質量則より遮音が悪くなります。二重窓は、2枚のガラスの間の空気がばねのはたらきをするため、ある低い振動数で共振して、かえって音を通しやすくなることもあります（波-2）。遮音の設計は、こうした効果を測定で確かめながら行います。

%% ---------------------------------------------------------------
\bksubsection{10}{響きをはかる：吸音と残響時間}

\leavevmode\label{ph:波-9:sec10}がらんとした体育館で手をたたくと、「ワーン」と長く響きます。カーテンやじゅうたんのある部屋では、すぐに消えます。音を出すのをやめたあと、部屋の中で音が残る長さを、数で表しましょう。

\noindent\textbf{吸音率と吸音力}\par\nobreak
音が壁に当たると、一部ははね返り、一部は壁にしみこんで熱になったり、壁を通りぬけたりします。はね返らなかった強さの割合を\textbf{吸音率} $\alpha$ といいます。\bkidx{きゅうおんりつ}{吸音率}コンクリートの壁は約0.02〜0.04（96〜98\,\% をはね返す）で、厚いカーテンや、音を吸うために作られた吸音材は、それよりずっと大きな値になります。壁の面積と吸音率の積を、部屋全体で足したもの $A=\alpha_1S_1+\alpha_2S_2+\cdots$〔m$^2$〕を\textbf{吸音力}といいます。「部屋に、音を全部吸いこむ窓が何 m$^2$ あいているのと同じか」と読めます。

\noindent\textbf{部屋のエネルギーのバランス式}\par\nobreak
体積 $V$ の部屋の中の音のエネルギーを $E$ とします。音はあらゆる向きに飛びかい、壁に当たるたびに、一部が吸いこまれます。詳しく数えると、1秒に壁に当たる音のエネルギーは、\qty{1}{m^2}あたり $cE/(4V)$ になります（いろいろな向きの音を平均した結果。4は向きの平均から出てくる数）。そのうち、吸音力 $A$ の分が吸いこまれるので、音を止めたあとは（たまる量＝出ていく量のマイナス。力-13 のバランス式）
\[ \frac{\dd E}{\dd t}=-\frac{cA}{4V}\,E \]
「減る速さが、そのときの量に比例する」形なので（6-4①）、エネルギーは指数関数で減っていきます。
\[ E=E_0\,e^{-cAt/(4V)} \]
エネルギーが100万分の1（60\,dB 下がる）になるまでの時間を\textbf{残響時間}といい、$t_{\mathrm R}$ と書きます。\bkidx{ざんきょうじかん}{残響時間}$e^{-cAt_{\mathrm R}/(4V)}=10^{-6}$ から、$cAt_{\mathrm R}/(4V)=6\ln10$ で
\[ t_{\mathrm R}=\frac{24\ln10}{c}\cdot\frac{V}{A} \]
です。$c=\qty{343}{m/s}$ を入れると、$24\times2.303\div343\approx0.161$ なので
\pheq{t_{\mathrm R}\approx0.161\,\frac{V}{A}}{s ＝ s/m × m$^3$÷m$^2$}
\label{ph:波-9:sabine}となります。これを\textbf{セイビンの式}といいます。\bkidx{せいびんのしき}{セイビンの式}0.161 の単位は s/m です。

\begin{phyomi}[$t_{\mathrm R}=0.161\,V/A$]
「部屋が大きいほど長く響き、よく吸う物が多いほど短く響く」。大きな部屋では、音が壁から壁まで飛ぶのに時間がかかり、壁に当たる回数が少ないので、なかなか減りません。縦\qty{7}{m}・横\qty{9}{m}・高さ\qty{3}{m}の教室で、吸音力が\qty{30}{m^2}なら、約1.0秒です（\bkreidairef{14}）。
\end{phyomi}

\begin{phbaisuru}{吸音力 $A$}
残響時間は半分になります。教室の天井（\qty{63}{m^2}、今の吸音率0.02）に吸音率0.7の吸音板をはると、吸音力はおよそ2.4倍になり、残響時間は約0.42秒になります。部屋の体積を2倍にすると（吸音力が同じなら）、残響時間も2倍です。
\par\vspace{2pt}\centering
\begin{tikzpicture}
\begin{axis}[bkplot, width=100mm, height=46mm, xmin=0, xmax=1.25, ymin=-66, ymax=4, xtick={0.25,0.5,0.75,1.0}, ytick={-60,-45,-30,-15,0}, xlabel={音を止めてからの時間（s）}, ylabel={レベル（dB）}, axis x line*=bottom, axis y line*=left, every axis x label/.style={at={(axis description cs:0.5,-0.17)}, anchor=north, font=\sffamily\small}]
  \addplot[bkcurve, domain=0:1.0] {-60*x/1.0};
  \addplot[bkOrange, line width=1.5pt, dashed, domain=0:0.5] {-60*x/0.5};
  \draw[bkGray, line width=0.5pt, densely dotted] (axis cs:0,-60) -- (axis cs:1.2,-60);
  \node[rxj lab, anchor=south west, text=bkBlue!80!black] at (axis cs:0.62,-30) {実線：$A=30$ m$^2$（1.0 s）};
  \node[rxj lab, anchor=south west, text=bkOrange!80!black] at (axis cs:0.30,-8) {破線：$A=60$ m$^2$（0.5 s）};
\end{axis}
\end{tikzpicture}
\end{phbaisuru}

dB で見ると、指数関数の減り方は直線になります（3-12⑥の片対数グラフと同じ）。残響時間をはかるときは、大きな音を止めてから、レベルが直線的に下がっていく傾きを読みとり、60\,dB 下がる時間に直します。

\begin{column}{響きの科学を始めたセイビン}
1895年ごろ、アメリカのハーバード大学で、新しい講堂の「話し声が聞きとりにくい」という困りごとが、若い物理学者セイビンにまかされました。セイビンは、夜の静かな時間に、オルガンのパイプの音を止めてから聞こえなくなるまでの時間を、ストップウォッチで何百回もはかりました。近くの劇場から座ぶとんを何百枚も運びこみ、その数を変えてはかったところ、「残響時間×（部屋の吸音力）」が、部屋ごとに決まった一定の値になることを見つけました。これがセイビンの式です。セイビンは、この式を使って、ボストンのコンサートホールの設計にも加わりました。響きを「計算して作る」建築音響の学問は、ここから始まったのです。
\end{column}

%% ---------------- たいせつ・ことば ----------------
\begin{taisetsu}
\begin{itemize}[leftmargin=1.2\zw, itemsep=2pt, topsep=0pt, label={\color{bkNavy}●}]
\item 音は圧力のゆれ（縦波）。音圧 $p=\rho_0cu$（$\rho_0c\approx\qty{413}{Pa.s/m}$）。真空では伝わらない。
\item 音の速さ $c=\sqrt{K/\rho_0}$。気体は $\sqrt{\gamma R_{\mathrm s}T}$、空気で $c\approx331.5+0.6\theta$。水 約\qty{1480}{m/s}、鋼の棒 約\qty{5100}{m/s}。
\item $L_p=20\log_{10}(p/p_0)$（$p_0=\qty{20}{\micro Pa}$）。音を合わせるときは強さをたす。同じ音2つで $+3$\,dB。
\item 点の音源：距離2倍で $-6$\,dB。線の音源：$-3$\,dB。$L_p\approx L_W-20\log_{10}R-11$。
\item 質量則 $\mathrm{TL}\approx20\log_{10}(\pi fm/(\rho_0c))$：重さ・振動数2倍で $+6$\,dB。残響時間 $t_{\mathrm R}\approx0.161V/A$。
\end{itemize}
\end{taisetsu}

\begin{kotoba}\small
\bkword{音圧}{音が来たときの、大気圧からの圧力のずれ。ふつうは実効値で表す。}
\bkword{音の強さ}{1秒に\qty{1}{m^2}を通る音のエネルギー（W/m$^2$）。音圧の2乗に比例する。}
\bkword{音圧レベル}{音圧を基準\qty{20}{\micro Pa}と比べて dB で表したもの（$20\log_{10}(p/p_0)$）。}
\bkword{A 特性}{耳の感じ方に合わせて、低い音などを弱めてから合わせる騒音計の補正。}
\bkword{透過損失・質量則}{壁を通る音がどれだけ小さくなるかを dB で表したもの。重い壁ほど大きい。}
\bkword{吸音力・残響時間}{部屋の音を吸う力（m$^2$）。音を止めてから60\,dB 下がるまでの時間。}
\end{kotoba}
\bkkigocards{\bkkigocard{$L_p$}{エル・ピー}{音圧レベル〔dB〕。基準は\qty{20}{\micro Pa}。音の強さのレベル $L_I$、音響パワーレベル $L_W$ も同じ仲間}{sound pressure level}{波-9}}{\bkkigocard{$t_{\mathrm R}$}{ティー・アール}{残響時間〔s〕。音を止めてからレベルが60\,dB 下がるまでの時間。専門書では $T$・$T_{60}$。波-9 では絶対温度 $T$ と分けた}{reverberation time}{波-9}}

\begin{senmonsho}[建築音響の本の書き方]
\begin{bkhonbun}
音圧レベルは $L_p=10\log_{10}(\overline{p^2}/p_0^{\,2})$、$p_0=\qty{20}{\micro Pa}$ で定義される。室内の残響時間は Sabine の式 $T=0.161\,V/A$、吸音力 $A=\sum S_i\alpha_i$ で与えられ、単層壁の音響透過損失は質量則 $R\approx20\log_{10}(fm)-42.5$〔dB〕で近似される。
\end{bkhonbun}
\bkyaku{$\overline{p^2}$ は「音圧の2乗の平均」で、実効値の2乗のこと。残響時間を $T$、透過損失を $R$ と書く本が多い（この本では $t_{\mathrm R}$ と TL。絶対温度 $T$・気体定数 $R$ と分けた）。$20\log_{10}(fm)-42.5$ は、⑨の $20\log_{10}(\pi fm/(\rho_0c))$ を、$20\log_{10}(\pi/413)\approx-42.4$ として書きかえた形。}
\end{senmonsho}

%% ---------------- 例題 ----------------
\Needspace{14\baselineskip}
\bksection{reidai}{例題でたしかめよう}

\noindent 空気の音の速さは、\qty{20}{\celsius}の\qty{343}{m/s}を使います。中学理科と比べる例題（雷・花火）は\qty{340}{m/s}とします。空気の $\rho_0c$ は\qty{413}{Pa.s/m}です。

\begin{phreidai}[雷までの距離]
\label{ph:波-9:ex1}稲妻が光ってから3.0秒後に、雷の音が聞こえました。雷までの距離は約何 km ですか。音の速さは\qty{340}{m/s}とします。
\phbamen
\bkhatenaref の場面です。光と音は、同じ所から同時に出ます。
\phmushi
光が届くまでの時間（\qty{1}{km}で100万分の3秒）。風。稲妻の長さ（光った所と音の出た所は同じとする）。
\phzu
雷までの距離 $d$、音の速さ $c=\qty{340}{m/s}$、おくれ $t=\qty{3.0}{s}$。
\phriyu
光の時間は無視できるので、おくれは「音が進むのにかかった時間」。速さ一定なので、距離＝速さ×時間（②）。
\phtateru
\[ d=c\,t \]
\phtoku
\[ d=340\times3.0=\qty{1020}{m}\approx\qty{1.0}{km} \]
\phtashikame
単位：m/s × s ＝ m。大きさ：「秒数÷3で km」の目安と合う。極端な場合：光ってすぐ音がしたら（$t\to0$）、雷はすぐ近く。
\kotae 約\qty{1.0}{km}
\end{phreidai}
\begin{jibun}
打ち上げ花火が開いてから2.5秒後に、ドーンという音が聞こえました。花火までの距離は約何 km ですか。音の速さは\qty{340}{m/s}とします。
\bkjans{$340\times2.5=\qty{850}{m}$、約\qty{0.85}{km}}
\end{jibun}

\begin{phreidai}[こだまで音の速さをはかる]
\label{ph:波-9:ex2}やってみようで、大きな壁から\qty{50}{m}離れて手をたたきました。こだまが聞こえるまでを10回続けてはかると、10回で2.9秒でした。音の速さを求めましょう。
\phbamen
手をたたいた音が、壁ではね返ってもどってくる。
\phmushi
手をたたく・時計を押す人のおくれ（10回続けて、ならしている）。風。壁は平らで、音はまっすぐ行ってもどる。
\phzu
壁までの距離 $d=\qty{50}{m}$、こだまの時間 $t$（1回分）、音の速さ $c$。
\phriyu
音は壁まで行って、もどってくるので、進んだ道のりは $2d$（往復）。距離＝速さ×時間。
\phtateru
\[ 2d=c\,t\qquad\text{より}\qquad c=\frac{2d}{t} \]
\phtoku
1回分の時間は $t=2.9\div10=\qty{0.29}{s}$。
\[ c=\frac{2\times50}{0.29}\approx\qty{345}{m/s} \]
\phtashikame
単位：m ÷ s ＝ m/s。大きさ：\qty{20}{\celsius}の\qty{343}{m/s}に近い（$331.5+0.6\theta=345$ なら約\qty{22}{\celsius}の日）。往復を忘れて $50\div0.29$ とすると約\qty{172}{m/s}になり、半分の速さでおかしい。
\kotae 約\qty{345}{m/s}
\end{phreidai}
\begin{jibun}
壁から\qty{85}{m}離れた所で、こだまが聞こえるまでの時間が0.50秒でした。音の速さを求めましょう。
\bkjans{往復で $2\times85=\qty{170}{m}$。$170\div0.50=\qty{340}{m/s}$}
\end{jibun}

\begin{reidai}[1 Pa の音で、空気はどれだけ動く？]
\label{ph:波-9:ex3}音圧\qty{1}{Pa}（実効値）の\qty{1000}{Hz}の音が、空気の中を進んでいます。空気が動く速さ（実効値）と、動く幅（ずれの実効値）を求めましょう。
\par\bkstep{理想化} 平面の波で、空気は正弦波の形でゆれるとする。
\par\kangaekata 速さは①の $u=p/(\rho_0c)$。正弦波のゆれでは、速さはずれの $2\pi f$ 倍（波-1）なので、ずれ＝速さ÷$(2\pi f)$。
\par\shiki $u=1\div413\approx\qty{2.4e-3}{m/s}$、$\xi=u\div(2\pi\times1000)\approx\qty{3.9e-7}{m}$
\par\bkstep{確かめ} ずれは約\qty{0.4}{\micro m}（\qty{1}{mm}の2500分の1）。94\,dB の大きな音でも、空気はわずかしか動かない。
\par\kotae 速さ約\qty{2.4}{mm/s}、ずれ約\qty{0.39}{\micro m}
\end{reidai}
\begin{jibun}
聞こえ始めの音圧\qty{20}{\micro Pa}の、\qty{1000}{Hz}の音では、空気が動く速さとずれ（実効値）はいくらですか。
\bkjans{速さは $20\times10^{-6}\div413$ $\approx\qty{4.8e-8}{m/s}$。ずれは $4.8\times10^{-8}\div6283$ $\approx\qty{7.7e-12}{m}$（原子の大きさの約30分の1）}
\end{jibun}

\begin{phreidai}[水の中の音の速さ]
\label{ph:波-9:ex4}\qty{20}{\celsius}の水の体積弾性率は約\qty{2.2}{GPa}、密度は\qty{998}{kg/m^3}です。水の中の音の速さを求めましょう。また、空気の中の何倍ですか。
\phbamen
プールにもぐって、遠くで石と石をたたき合わせた音を聞く。
\phmushi
水の温度・圧力による $K$ と $\rho_0$ の変わり方。水にとけた空気の泡。
\phzu
体積弾性率 $K=\qty{2.2e9}{Pa}$、密度 $\rho_0=\qty{998}{kg/m^3}$、音の速さ $c$。
\phriyu
水は液体なので、縮みにくさは体積弾性率 $K$ で、③の $c=\sqrt{K/\rho_0}$ が使える。音のような速い変化でも、水は温度がほとんど変わらないので、$K$ はそのままでよい。
\phtateru
\[ c=\sqrt{\frac{K}{\rho_0}} \]
\phtoku
\[ c=\sqrt{\frac{2.2\times10^9}{998}}=\sqrt{2.20\times10^6}\approx\qty{1480}{m/s},\qquad 1480\div343\approx4.3 \]
\phtashikame
単位：$\sqrt{\text{Pa}\div(\text{kg/m}^3)}=\sqrt{(\text{kg/(m·s}^2))\times\text{m}^3/\text{kg}}$＝m/s。大きさ：はかった値（約\qty{1480}{m/s}）と合う。極端な場合：$K$ が空気と同じくらいなら、重い水の中の音はずっと遅くなるはず。水がとても縮みにくいので速い。
\kotae 約\qty{1480}{m/s}、空気の約4.3倍
\end{phreidai}
\begin{jibun}
海水の体積弾性率を約\qty{2.34}{GPa}、密度を\qty{1025}{kg/m^3}として、海の中の音の速さを求めましょう。
\bkjans{$\sqrt{2.34\times10^9\div1025}\approx\qty{1510}{m/s}$（約\qty{1500}{m/s}）}
\end{jibun}

\begin{phreidai}[レールを伝わる音]
\label{ph:波-9:ex5}\qty{1.0}{km}先で、鋼のレールをハンマーでたたきました。レールを伝わる音と、空気を伝わる音は、それぞれ何秒後に届きますか。鋼は $E=\qty{205}{GPa}$、$\rho_0=\qty{7850}{kg/m^3}$、空気の音の速さは\qty{343}{m/s}とします。
\phbamen
線路の点検で、遠くでレールをたたいた音を、レールに当てた器械と、耳の両方で聞く。
\phmushi
レールの中で音が弱まること。レールの太さ（波長より細い棒とみる）。
\phzu
距離 $d=\qty{1000}{m}$、鋼の棒の音の速さ $c_{\text{鋼}}$、空気の $c_{\text{空気}}=\qty{343}{m/s}$。
\phriyu
レールは細長い棒なので、縮みにくさはヤング率 $E$ で、$c=\sqrt{E/\rho_0}$（④、波-5 ⑤）。時間＝距離÷速さ。
\phtateru
\[ c_{\text{鋼}}=\sqrt{\frac{E}{\rho_0}},\qquad t=\frac{d}{c} \]
\phtoku
\[ c_{\text{鋼}}=\sqrt{\frac{205\times10^9}{7850}}\approx\qty{5110}{m/s} \]
\[ t_{\text{鋼}}=\frac{1000}{5110}\approx\qty{0.20}{s},\qquad t_{\text{空気}}=\frac{1000}{343}\approx\qty{2.9}{s} \]
\phtashikame
単位：s。大きさ：鋼の速さは力-16 の約\qty{5100}{m/s}と合う。空気の約15倍速いので、時間は約15分の1。
\kotae レールは約0.20秒後、空気は約2.9秒後
\end{phreidai}
\begin{jibun}
\qty{500}{m}先でレールをたたいたとき、レールを伝わる音（\qty{5100}{m/s}）と、空気を伝わる音（\qty{343}{m/s}）が届く時間の差は何秒ですか。
\bkjans{$500\div343\approx1.458$、$500\div5100\approx0.098$。差は約1.36秒}
\end{jibun}

\begin{reidai}[\qty{15}{\celsius}の音の速さを2通りで]
\label{ph:波-9:ex6}\qty{15}{\celsius}の空気の音の速さを、(1) $331.5+0.6\theta$ と、(2) $\sqrt{\gamma R_{\mathrm s}T}$（$\gamma=1.4$、$R_{\mathrm s}=\qty{287}{J/(kg.K)}$）で求めて、比べましょう。
\par\bkstep{理想化} 乾いた空気を理想気体とする（湿度による小さな違いを無視）。
\par\kangaekata (1) は⑤の1次式。(2) は $T=273.15+15=\qty{288.15}{K}$ を入れる。
\par\shiki (1) $331.5+0.6\times15=340.5$
\par\hspace*{2\zw}(2) $\sqrt{1.4\times287\times288.15}=\sqrt{\num{115779}}\approx340.3$
\par\bkstep{確かめ} 差は\qty{0.2}{m/s}で、1次近似はよく合う。中学理科の「約\qty{340}{m/s}」は\qty{15}{\celsius}くらいの値。
\par\kotae (1) \qty{340.5}{m/s}、(2) 約\qty{340.3}{m/s}
\end{reidai}
\begin{jibun}
\qty{-10}{\celsius}の冬の朝の、空気の音の速さを (1) $331.5+0.6\theta$ と (2) $\sqrt{\gamma R_{\mathrm s}T}$ で求めましょう。
\bkjans{(1) $331.5-6=\qty{325.5}{m/s}$\quad (2) $\sqrt{1.4\times287\times263.15}\approx\qty{325.2}{m/s}$}
\end{jibun}

\begin{reidai}[1 Pa は何 dB？]
\label{ph:波-9:ex7}音圧が\qty{1}{Pa}の音の音圧レベルを求めましょう。$\log_{10}2=0.301$ を使います。
\par\bkstep{理想化} 音圧は実効値とする。
\par\kangaekata ⑥の $L_p=20\log_{10}(p/p_0)$、$p_0=\qty{20e-6}{Pa}$。
\par\shiki $p/p_0=1\div(20\times10^{-6})=5\times10^4$
\par\hspace*{2\zw}$\log_{10}(5\times10^4)=4+\log_{10}5=4+(1-0.301)=4.699$
\par\hspace*{2\zw}$L_p=20\times4.699\approx94.0$
\par\bkstep{確かめ} 会話の\qty{0.02}{Pa}（60\,dB）の50倍で、$20\log_{10}50\approx34$\,dB 大きい。$60+34=94$ と合う。
\par\kotae 約94\,dB
\end{reidai}
\begin{jibun}
(1) 音圧\qty{0.2}{Pa}は何 dB ですか。(2) 100\,dB の音の音圧は何 Pa ですか。
\bkjans{(1) $20\log_{10}(0.2\div0.00002)=20\log_{10}10^4=80$\,dB\quad (2) $100\div20=5$ なので $p=20\times10^{-6}\times10^5=\qty{2}{Pa}$}
\end{jibun}

\begin{reidai}[会話とささやき声の音圧の比]
\label{ph:波-9:ex8}ふつうの会話は60\,dB、ささやき声は30\,dB です（3-12 と同じ数）。音圧は何倍違いますか。強さは何倍ですか。$\sqrt{10}\approx3.16$ を使います。
\par\bkstep{理想化} どちらも空気の中の音で、dB は音圧レベルと音の強さのレベルで同じとする（⑥）。
\par\kangaekata 差は30\,dB。音圧は「20\,dB で10倍」、強さは「10\,dB で10倍」。
\par\shiki 音圧：$10^{30/20}=10^{1.5}=10\times\sqrt{10}\approx31.6$。強さ：$10^{30/10}=10^3=1000$
\par\bkstep{確かめ} 音圧の比の2乗が強さの比（$31.6^2\approx1000$）になっている。
\par\kotae 音圧は約32倍、強さは1000倍
\end{reidai}
\begin{jibun}
工場の機械の近く（90\,dB）と、静かな図書館（40\,dB）では、音圧と強さは何倍違いますか（3-12 と同じ数）。
\bkjans{差は50\,dB。音圧 $10^{2.5}\approx316$ 倍、強さ $10^5=\num{100000}$ 倍}
\end{jibun}

\begin{phreidai}[機械の音を合わせる]
\label{ph:波-9:ex9}ある場所で、機械 P を1台だけ動かすと80\,dB でした。(1) 同じ機械をもう1台、同じ場所に置いて同時に動かすと、何 dB になりますか。(2) 1台の機械 P（80\,dB）と、70\,dB の機械 Q を同時に動かすと、何 dB になりますか。$\log_{10}2=0.301$、$\log_{10}1.1=0.041$ を使います。
\phbamen
工場に機械を増やしたときの騒音を見積もる。
\phmushi
背景の音（周りの音）。2台の音の山と谷がそろうこと（互いに関係なくゆれているとする）。
\phzu
レベル $L_1=80$\,dB、$L_2$、合わせたレベル $L$。基準の強さとの比 $10^{L/10}$。
\phriyu
関係なくゆれる2つの音は、強さがたし算になる（⑦）。dB はそのままたせないので、強さの比にもどしてからたす。
\phtateru
\[ L=10\log_{10}\bigl(10^{L_1/10}+10^{L_2/10}\bigr) \]
\phtoku
(1) $10^{8}+10^{8}=2\times10^8$ なので
\[ L=10\log_{10}(2\times10^8)=80+10\log_{10}2=80+3.01\approx\qty{83}{dB} \]
(2) $10^{8}+10^{7}=1.1\times10^8$ なので
\[ L=10\log_{10}(1.1\times10^8)=80+10\times0.041\approx\qty{80.4}{dB} \]
\phtashikame
大きさ：合わせると、大きいほうより少しだけ大きくなる（$80+80=160$ のようにはならない）。(2) は、差が10\,dB なので表\ref{ph:波-9:tab-tashi}の0.4\,dB と合う。極端な場合：Q がとても小さい音なら、ほぼ80\,dB のまま。
\kotae (1) 約83\,dB\quad (2) 約80.4\,dB
\end{phreidai}
\begin{jibun}
(1) 70\,dB の機械を3台同時に動かすと、何 dB になりますか（$\log_{10}3=0.477$）。(2) 90\,dB と80\,dB の音を合わせると、何 dB ですか。
\bkjans{(1) $70+10\log_{10}3=70+4.77\approx\qty{74.8}{dB}$\quad (2) $90+10\log_{10}1.1\approx\qty{90.4}{dB}$}
\end{jibun}

\begin{phreidai}[スピーカーから離れた所の音]
\label{ph:波-9:ex10}音響パワーレベル100\,dB の小さなスピーカーが、周りに何もない広い空中に置かれています。\qty{10}{m}離れた所の音圧レベルを求めましょう。スピーカーの音響パワーは何 W ですか。
\phbamen
野外のイベントで、スピーカーから離れた客席の音の大きさを見積もる。
\phmushi
地面や建物ではね返る音。空気が音を吸う分。スピーカーの向き（どの向きにも同じ強さで出すとする）。
\phzu
音響パワーレベル $L_W=100$\,dB、距離 $R=\qty{10}{m}$、音圧レベル $L_p$。
\phriyu
小さな音源から四方へ広がる音は、半径 $R$ の球面に広がる（点の音源。波-7 ②）。⑦の $L_p\approx L_W-20\log_{10}R-11$。
\phtateru
\[ L_p=L_W-20\log_{10}R-11,\qquad W=W_0\times10^{L_W/10} \]
\phtoku
\[ L_p=100-20\log_{10}10-11=100-20-11=\qty{69}{dB} \]
\[ W=10^{-12}\times10^{10}=\qty{0.01}{W} \]
\phtashikame
単位：dB。大きさ：\qty{10}{m}で69\,dB は、にぎやかな通りくらい。0.01 W の音でも、かなり大きく聞こえる（スピーカーは、電気の何\,\% かしか音にならない）。極端な場合：\qty{20}{m}なら6\,dB 下がって63\,dB。
\kotae 約69\,dB、\qty{0.01}{W}
\end{phreidai}
\begin{jibun}
音響パワーレベル95\,dB の点の音源から、\qty{20}{m}の所の音圧レベルを求めましょう（周りに何もないとします。$\log_{10}2=0.301$）。
\bkjans{$95-20\log_{10}20-11=95-26.0-11\approx\qty{58}{dB}$}
\end{jibun}

\begin{reidai}[道路の音と工事の音]
\label{ph:波-9:ex11}\qty{10}{m}離れた所で、どちらも70\,dB の、まっすぐな道路の音（線の音源）と、工事の機械の音（点の音源）があります。\qty{40}{m}離れると、それぞれ何 dB になりますか。
\par\bkstep{理想化} 周りに何もなく、音ははね返らないとする。道路は長いまっすぐな線とする。
\par\kangaekata 距離が4倍＝2倍の2回。点の音源は2倍ごとに $-6$\,dB、線の音源は $-3$\,dB（⑦）。
\par\shiki 点：$70-6\times2=58$。線：$70-3\times2=64$
\par\bkstep{確かめ} 式で：点 $70-20\log_{10}4\approx58.0$、線 $70-10\log_{10}4\approx64.0$ で合う。線の音源のほうが、離れても小さくなりにくい。
\par\kotae 道路は約64\,dB、工事の機械は約58\,dB
\end{reidai}
\begin{jibun}
\qty{20}{m}の所で75\,dB の道路の音は、\qty{80}{m}の所では何 dB ですか（線の音源とします）。
\bkjans{距離4倍で $-3\times2=-6$\,dB。約69\,dB}
\end{jibun}

\begin{phreidai}[低い音は小さく聞こえる：A 特性]
\label{ph:波-9:ex12}音圧レベルがどれも80\,dB の、\qty{100}{Hz}・\qty{1000}{Hz}・\qty{4000}{Hz}の音があります。A 特性の補正は、それぞれ $-19.1$\,dB、0\,dB、$+1.0$\,dB です。それぞれ何 dB(A) ですか。
\phbamen
騒音計で、低い音・中くらいの音・高い音をはかる。
\phmushi
人による聞こえ方の違い（A 特性は、多くの人の平均に近づけた決まり）。
\phzu
音圧レベル $L_p=80$\,dB、A 特性の補正 $\Delta L_{\mathrm A}$、A 特性音圧レベル $L_{\mathrm A}$。
\phriyu
A 特性は、振動数ごとに、補正の dB をたすと決められている（⑧）。1つの高さの音なら、たし算だけでよい。
\phtateru
\[ L_{\mathrm A}=L_p+\Delta L_{\mathrm A} \]
\phtoku
\qty{100}{Hz}：$80-19.1=60.9$。\qty{1000}{Hz}：$80+0=80$。\qty{4000}{Hz}：$80+1.0=81$
\phtashikame
大きさ：低い音は、同じ音圧でも、ずっと小さな dB(A) になる（耳が鈍い）。極端な場合：\qty{1000}{Hz}では補正が0なので、dB と dB(A) は同じ。
\kotae \qty{100}{Hz}は約61\,dB(A)、\qty{1000}{Hz}は80\,dB(A)、\qty{4000}{Hz}は約81\,dB(A)
\end{phreidai}
\begin{jibun}
\qty{63}{Hz}で90\,dB の低い音（A 特性の補正 $-26.2$\,dB）と、\qty{1000}{Hz}で70\,dB の音が、同時に鳴っています。(1) 補正なしの合計（dB）と、(2) A 特性の合計（dB(A)）を求めましょう。
\bkjans{(1) $10\log_{10}(10^9+10^7)\approx\qty{90.0}{dB}$\quad (2) $90-26.2=63.8$ と70 を合わせて、$10\log_{10}(10^{6.38}+10^7)\approx\qty{70.9}{dB}$(A)}
\end{jibun}

\begin{phreidai}[窓ガラスの遮音]
\label{ph:波-9:ex13}厚さ\qty{5}{mm}のガラス（面密度 $m=\qty{12.5}{kg/m^2}$）の、\qty{500}{Hz}の音に対する透過損失を、垂直に当たる音の質量則で求めましょう。外の道路の音が\qty{500}{Hz}で75\,dB なら、部屋の中ではおよそ何 dB ですか。
\phbamen
道路に面した部屋の窓で、外の音をどれだけ小さくできるかを見積もる。
\phmushi
窓のすきま、窓わくを通る音。ガラスのしなり（板のこわさ）と、その中でのむだ（質量則は「重さだけの板」）。部屋の中での響き。
\phzu
面密度 $m$、振動数 $f$、空気の $\rho_0c=\qty{413}{Pa.s/m}$、透過損失 TL。
\phriyu
重さだけで決まる遮音なので、⑨の質量則。窓の外と中のレベルの差が透過損失。
\phtateru
\[ \mathrm{TL}\approx20\log_{10}\frac{\pi fm}{\rho_0c} \]
\phtoku
\[ \frac{\pi fm}{\rho_0c}=\frac{3.14\times500\times12.5}{413}\approx47.5 \]
\[ \mathrm{TL}\approx20\log_{10}47.5\approx20\times1.677\approx\qty{34}{dB} \]
部屋の中は $75-34=\qty{41}{dB}$。
\phtashikame
単位：$\pi fm/(\rho_0c)$ は（1/s × kg/m$^2$）÷（Pa·s/m）で単位なし。大きさ：ガラス1枚で30\,dB 台は、カタログの値と同じくらい（本当は斜めの音もあり、もう少し小さい：⑨のせいかくに言うと）。極端な場合：ガラスを2倍の厚さにすると、40\,dB になる。
\kotae 約34\,dB、部屋の中は約41\,dB
\end{phreidai}
\begin{jibun}
厚さ\qty{1}{mm}の鉄板（$m=\qty{7.85}{kg/m^2}$）の、\qty{1000}{Hz}の音に対する透過損失を、質量則で求めましょう。
\bkjans{$\pi\times1000\times7.85\div413\approx59.7$。$20\log_{10}59.7\approx\qty{36}{dB}$}
\end{jibun}

\begin{phreidai}[教室の残響時間]
\label{ph:波-9:ex14}縦\qty{7}{m}・横\qty{9}{m}・高さ\qty{3}{m}の教室の、吸音力が\qty{30}{m^2}です。残響時間を求めましょう。また、壁や天井などの表面の吸音率の平均はいくらですか。
\phbamen
先生の声が響きすぎて聞きとりにくい教室を、数で調べる。
\phmushi
机やいす・人による吸音の場所のかたより（音は部屋じゅうに一様に広がるとする）。空気が音を吸う分。
\phzu
体積 $V$、表面積 $S$、吸音力 $A=\qty{30}{m^2}$、残響時間 $t_{\mathrm R}$、平均の吸音率 $\bar\alpha=A/S$。
\phriyu
音を止めてからの減り方は、部屋のエネルギーのバランス式で決まり、⑩のセイビンの式になる。
\phtateru
\[ t_{\mathrm R}=0.161\,\frac{V}{A},\qquad \bar\alpha=\frac{A}{S} \]
\phtoku
$V=7\times9\times3=\qty{189}{m^3}$。$S=2\times(7\times9)+2\times(7+9)\times3=126+96=\qty{222}{m^2}$。
\[ t_{\mathrm R}=0.161\times\frac{189}{30}\approx\qty{1.0}{s},\qquad \bar\alpha=\frac{30}{222}\approx0.14 \]
\phtashikame
単位：s/m × m$^3$÷m$^2$＝s。大きさ：話を聞く部屋としては少し長い（話す部屋は1秒より短めがよいとされる）。極端な場合：吸音力がとても小さい（$A\to0$）と、いつまでも響く。
\kotae 約1.0秒、平均の吸音率は約0.14
\end{phreidai}
\begin{jibun}
\bkreidairef{14}の教室で、吸音力を\qty{60}{m^2}にすると、残響時間は何秒になりますか。
\bkjans{$0.161\times189\div60\approx\qty{0.51}{s}$（半分）}
\end{jibun}

\begin{reidai}[聞こえる音の波長]
\label{ph:波-9:ex15}人に聞こえる音は、およそ\qty{20}{Hz}から\qty{20000}{Hz}までです。空気の中（\qty{343}{m/s}）での、それぞれの波長を求めましょう。
\par\bkstep{理想化} 音の速さは振動数によらない（空気の中では、ほぼ成り立つ）。
\par\kangaekata $c=f\lambda$（波-5 ②）から $\lambda=c/f$。
\par\shiki $343\div20\approx17.2$、$343\div\num{20000}\approx0.0172$
\par\bkstep{確かめ} 1000倍違う振動数で、波長も1000倍違う。\qty{1000}{Hz}では約\qty{0.34}{m}。
\par\kotae 約\qty{17}{m}から約\qty{17}{mm}まで
\end{reidai}
\begin{jibun}
水の中（\qty{1480}{m/s}）では、\qty{20}{Hz}と\qty{20000}{Hz}の音の波長は、それぞれいくらですか。
\bkjans{$1480\div20=\qty{74}{m}$、$1480\div\num{20000}=\qty{0.074}{m}$（\qty{74}{mm}）}
\end{jibun}

%% ---------------- シミュレーション ----------------
\Needspace{10\baselineskip}
\noindent\textbf{プログラム：騒音計の中でしていること}\par\nobreak
本当の騒音は、いろいろな高さの音がまざっています。騒音計は、音を高さごと（オクターブごと）に分け、それぞれに A 特性の補正をたしてから、強さをたし合わせます（⑦⑧）。換気扇の音を、高さごとにはかった値で計算してみましょう。A 特性の補正は、規格の式で計算します。
\begin{phcode}[Python]
import math
f = [63, 125, 250, 500, 1000, 2000, 4000, 8000]  \# Hz
L = [78, 75, 72, 70, 68, 65, 60, 55]              \# dB
def A(f):                       \# A 特性の補正
    f2 = f * f
    r = 12194**2 * f2**2 / ((f2 + 20.6**2) * (f2 + 12194**2)
        * math.sqrt((f2 + 107.7**2) * (f2 + 737.9**2)))
    return 20 * math.log10(r) + 2.00
def tasu(Ls):                   \# 強さにもどしてたす
    return 10 * math.log10(sum(10**(x / 10) for x in Ls))
print('補正なし', round(tasu(L), 1), 'dB')
LA = [x + A(fi) for x, fi in zip(L, f)]
print('A 特性 ', round(tasu(LA), 1), 'dB(A)')
\end{phcode}
\begin{phcodeout}
補正なし 81.2 dB
A 特性  73.0 dB(A)
\end{phcodeout}
この換気扇は、低い音がいちばん大きいので、A 特性でまとめると、補正なしより約8\,dB 小さくなります。耳に聞こえる「うるささ」は、補正なしの数ほどではないのです。

\noindent L の数を変えて、どの高さの音を小さくすると効くか、試してみましょう。いちばん大きい\qty{63}{Hz}の音を10\,dB 下げる（L の最初を68にする）と、補正なしは78.7\,dB まで下がりますが、A 特性は72.9\,dB(A) とほとんど変わりません。\qty{1000}{Hz}の音を10\,dB 下げると、補正なしは81.0\,dB とほぼそのままなのに、A 特性は71.5\,dB(A) に下がります。耳のうるささを減らすには、耳が敏感な高さの音を減らすのが効くのです。

%% ---------------- まちがえやすいところ ----------------
\Needspace{9\baselineskip}
\bksection{tsumazuki}{まちがえやすいところをチェック}
\begin{tsumazuki}{こだまの距離で、往復を忘れる}
\bkng{壁まで\qty{50}{m}、こだまが0.29秒なら、音の速さは $50\div0.29\approx\qty{172}{m/s}$。}
\bkok{音は壁まで行って帰るので、$2\times50\div0.29\approx\qty{345}{m/s}$。}
こだま・ソナー・魚群探知機など、はね返りを使う計算では、いつも「往復」です（波-10 でも使います）。
\end{tsumazuki}

\begin{tsumazuki}{dB をそのままたす}
\bkng{80\,dB の機械を2台動かすと、$80+80=160$\,dB。}
\bkok{強さが2倍になるので、$80+3=83$\,dB。}
dB をたしてよいのは、増幅器の倍率をかけるときです。音を合わせるときは、強さにもどしてたします。
\end{tsumazuki}

\begin{tsumazuki}{音圧の dB に10をかける}
\bkng{音圧が10倍になると、$10\log_{10}10=10$\,dB 大きくなる。}
\bkok{強さは音圧の2乗なので100倍。$20\log_{10}10=20$\,dB 大きくなる。}
強さ（エネルギー）の比には10、音圧の比には20をかけます。
\end{tsumazuki}

\begin{tsumazuki}{山の上では、気圧が低いから音が遅い？}
\bkng{富士山の頂上は気圧が低いので、音の速さも遅くなる。}
\bkok{気体の音の速さは $\sqrt{\gamma R_{\mathrm s}T}$ で、圧力によらない。頂上で遅いのは、気温が低いから。}
圧力が下がると、縮みにくさ $\gamma p_0$ も、密度 $\rho_0$ も同じ割合で小さくなり、打ち消し合います。
\end{tsumazuki}

\begin{tsumazuki}{距離が2倍で、dB も半分になる？}
\bkng{\qty{10}{m}で80\,dB の音は、\qty{20}{m}では40\,dB。}
\bkok{点の音源なら、強さが $1/4$ になって $-6$\,dB。\qty{20}{m}で74\,dB。}
dB は「けた」を数えた数なので、強さが半分でも dB は半分になりません。
\end{tsumazuki}

%% ---------------- 工学ではこう使う ----------------
\Needspace{8\baselineskip}
\bksection{phgear}{工学ではこう使う}
\begin{phkougaku}[騒音をはかる・防ぐ・響きをつくる]
\phboxhead{① 騒音計と、騒音の決まり}
騒音計は、マイクロホンで音圧をはかり、A 特性の補正をして、dB(A) で表示します。性能は規格で決められていて、正しくはかるには、使う前に、決まった大きさの音を出す器械で合わせます（校正）。日本では、住宅の地域の騒音について、国が「昼（6時〜22時）は55\,dB 以下、夜は45\,dB 以下が望ましい」という目安（騒音に係る環境基準。A 特性の等価騒音レベルで比べる）を決めています（幹線道路のすぐそばは、昼70\,dB 以下・夜65\,dB 以下）。工場や工事の現場では、85\,dB 以上の所で働く人の耳を守る対策（耳せん・作業時間の工夫）が、国のガイドラインですすめられています。

\begin{phketa}[音の数の目安]
\phketaline{空気・水・鋼（棒）の音の速さ}{\qty{343}{m/s}・約\qty{1480}{m/s}・約\qty{5100}{m/s}}
\phketaline{空気の $\rho_0c$}{約\qty{413}{Pa.s/m}}
\phketaline{会話・ガード下の音圧レベル}{約60\,dB（\qty{0.02}{Pa}）・約100\,dB（\qty{2}{Pa}）}
\phketaline{ガラス\qty{5}{mm}の透過損失（\qty{500}{Hz}）}{約30\,dB}
\phketaline{教室・コンサートホールの残響時間}{約0.6秒・約2秒前後}
\phboxhead{答えが桁で合っているか}
空気の中の音の速さが\qty{300}{m/s}台から外れたら、温度を K にしたか、往復を忘れていないかを見直す。身近な音の音圧レベルは0〜130\,dB に入る。dB の足し算で、合計が大きいほうより10\,dB 以上大きくなったら、dB をそのままたしていないかを見直す。
\end{phketa}

\phboxhead{② 道路の騒音を減らす}
高速道路の横に立つ遮音壁は、音の通り道をさえぎり、壁のかげになる所の音を小さくします。音は壁の上を回りこむ（回折。波-8）ので、壁は高いほど、また家に近いほど効きます。道路の表面に小さなすき間の多い舗装を使い、タイヤの音を吸わせる方法もあります。
\par\vspace{2pt}{\centering
\begin{tikzpicture}[x=1mm, y=1mm]
  \fill[bkGray!25] (0,0) rectangle (112,3);
  \draw[bkInk, line width=0.6pt] (0,3) -- (112,3);
  \filldraw[fill=bkBlue!35, draw=bkInk, line width=0.6pt, rounded corners=1pt] (6,3.5) rectangle (22,9) (9,9) rectangle (19,12.5);
  \fill[bkInk] (10,3.5) circle (1.4) (18,3.5) circle (1.4);
  \filldraw[fill=bkGray!55, draw=bkInk, line width=0.7pt] (44,3) rectangle (47,26);
  \node[rxj lab, anchor=west] at (48,16) {遮音壁};
  \filldraw[fill=bkOrange!25, draw=bkInk, line width=0.6pt] (92,3) rectangle (108,15);
  \filldraw[fill=bkRed!40, draw=bkInk, line width=0.6pt] (90,15) -- (100,22) -- (110,15) -- cycle;
  \node[rxj lab, anchor=north] at (100,-0.5) {家};
  \node[rxj lab, anchor=north] at (12,-0.5) {車（音源）};
  \draw[rxj arr, bkGray, densely dotted] (22,7) -- (43,9.5);
  \node[rxj lab, anchor=north, text=bkGray!80!black] at (66,-0.5) {点線：さえぎられる音};
  \draw[rxj arr, bkBlue!80!black] (22,9) -- (45,27) .. controls (60,24) and (80,16) .. (91,11);
  \node[rxj lab, anchor=south west, text=bkBlue!70!black] at (62,25.5) {実線：上を回りこむ音（小さくなる）};
\end{tikzpicture}\par}

\phboxhead{③ ホールの響きをつくる}
コンサートホールでは、音楽が豊かに聞こえるように、残響時間を約2秒前後にすることが多く、話を聞く講堂や教室では、言葉がはっきり聞こえるように短くします（アメリカの教室の規格では、ふつうの大きさの教室で0.6秒以下）。設計では、セイビンの式で吸音材の量を決め、模型やコンピューターで、音のはね返り方も確かめます。客席に人が座ると吸音力が大きく変わるので、いすにも人と同じくらい音を吸う布をはるくふうがあります。
\end{phkougaku}
\phstation{音響工学,建築音響,騒音・振動の対策,機械の静音設計}{音速、体積弾性率、音圧、実効値、音の強さ、音圧レベル、音響パワーレベル、逆2乗の法則、A 特性、等ラウドネス曲線、透過損失、質量則、コインシデンス、吸音率、吸音力、残響時間、セイビンの式}

\begin{phsaki}{音響工学・建築音響}
この章では、音を「小さな圧力のゆれ」として、平面の波と、部屋じゅうに一様に広がった音で考えました。小さな部屋の中では、音は定在波（波-6 ⑧）になり、場所によって大きく聞こえたり小さく聞こえたりします。部屋の音を波として正しく計算する方法や、壁の中の音の伝わり方、スピーカーやマイクロホンのしくみは、音響工学・建築音響の専門の本で学びます。爆発やジェットエンジンのような、とても大きな音では、音の波の形がくずれて衝撃波になります（非線形音響。熱-4 ⑨）。人の耳と脳が、どのように音を聞き分けているかは、聴覚の研究の分野で調べられています。
\end{phsaki}

%% ---------------- 章末問題 ----------------
\bkendproblems
\bkproblemgroup{基本}{3}
\bkproblem 打ち上げ花火が開いてから、4.0秒後に音が聞こえました。花火までの距離は約何 km ですか。音の速さは\qty{340}{m/s}とします（\nobreak\bkreidairef{1}\nobreak）。
\bkpans{$340\times4.0=\qty{1360}{m}$、約\qty{1.4}{km}}
\bkproblem \qty{25}{\celsius}の空気の音の速さを $331.5+0.6\theta$ で求め、\qty{440}{Hz}の音の波長を求めましょう（\nobreak\bkreidairef{6}・\bkreidairef{15}\nobreak）。
\bkpans{$331.5+0.6\times25=\qty{346.5}{m/s}$。$346.5\div440\approx\qty{0.79}{m}$}
\bkproblem 壁から\qty{68}{m}離れて手をたたいたら、こだまが0.40秒後に聞こえました。音の速さを求めましょう（\nobreak\bkreidairef{2}\nobreak）。
\bkpans{$2\times68\div0.40=\qty{340}{m/s}$}
\bkproblem (1) 音圧\qty{0.002}{Pa}は何 dB ですか。(2) 120\,dB の音の音圧は何 Pa ですか（\nobreak\bkreidairef{7}\nobreak）。
\bkpans{(1) $20\log_{10}(0.002\div0.00002)=20\log_{10}100=40$\,dB\quad (2) $20\times10^{-6}\times10^{6}=\qty{20}{Pa}$}
\bkproblem 85\,dB の機械が4台あります。全部同時に動かすと何 dB ですか。また、85\,dB と75\,dB の機械を合わせると何 dB ですか（\nobreak\bkreidairef{9}\nobreak）。
\bkpans{4台：$85+10\log_{10}4=85+6.0=\qty{91}{dB}$。85と75：$85+10\log_{10}1.1\approx\qty{85.4}{dB}$}
\bkproblem 音響パワーレベル110\,dB の機械を、周りに何もない所に置きました（点の音源）。\qty{30}{m}離れた所の音圧レベルを求めましょう（$\log_{10}3=0.477$）（\nobreak\bkreidairef{10}\nobreak）。
\bkpans{$110-20\log_{10}30-11=110-29.5-11\approx\qty{69.5}{dB}$（約70\,dB）}
\bkproblem 音圧が\qty{0.063}{Pa}の音の、音圧レベルと音の強さを求めましょう（$\log_{10}3.15=0.498$、空気の $\rho_0c=\qty{413}{Pa.s/m}$）（\nobreak\bkreidairef{7}\nobreak）。
\bkpans{$p/p_0=0.063\div0.00002=3150$。$20\times(3+0.498)\approx\qty{70}{dB}$。強さは $0.063^2\div413\approx\qty{9.6e-6}{W/m^2}$}
\bkproblem まっすぐな道路の音（線の音源）と、工事の機械の音（点の音源）が、どちらも\qty{20}{m}の所で72\,dB です。\qty{80}{m}の所では、それぞれ何 dB ですか。2つを合わせると何 dB ですか（表2 を使う）（\nobreak\bkreidairef{9}・\bkreidairef{11}\nobreak）。
\bkpans{距離4倍で、道路は $72-3\times2=\qty{66}{dB}$、機械は $72-6\times2=\qty{60}{dB}$。差が6\,dB なので表2 から1.0\,dB たして、約\qty{67}{dB}}
\bkproblem 冬（$-5$\,℃）と夏（35\,℃）に、稲妻が光って3.0秒後に雷の音が聞こえました。$331.5+0.6\theta$ で音の速さを求め、雷までの距離の差を求めましょう（\nobreak\bkreidairef{1}・\bkreidairef{6}\nobreak）。
\bkpans{冬 $328.5\times3.0\approx\qty{986}{m}$、夏 $352.5\times3.0\approx\qty{1058}{m}$。差は約\qty{72}{m}（「秒数÷3で km」の目安で十分）}
\bkproblemgroup{チャレンジ}{4}
\bkproblem 厚さ\qty{8}{mm}のガラス（$m=\qty{20}{kg/m^2}$）の、\qty{500}{Hz}の透過損失を質量則で求めましょう。\qty{5}{mm}のガラス（34\,dB）より何 dB 大きいですか（\nobreak\bkreidairef{13}\nobreak）。
\bkpans{$\pi\times500\times20\div413\approx76.1$、$20\log_{10}76.1\approx\qty{38}{dB}$。約4\,dB 大きい（$20\log_{10}(20/12.5)\approx4.1$）}
\bkproblem 体積\qty{6000}{m^3}、吸音力\qty{900}{m^2}の体育館の残響時間を求めましょう（\nobreak\bkreidairef{14}\nobreak）。
\bkpans{$0.161\times6000\div900\approx\qty{1.1}{s}$}
\bkproblem 遠くで鋼のレールをたたいたら、レールを伝わる音（\qty{5100}{m/s}）が聞こえてから1.0秒後に、空気を伝わる音（\qty{343}{m/s}）が聞こえました。たたいた所までの距離を求めましょう（\nobreak\bkreidairef{5}\nobreak）。
\bkpans{距離を $d$ とすると $d/343-d/5100=1.0$。$1/343-1/5100$ $\approx0.0029155-0.0001961$ $=0.0027194$。$d=1.0\div0.0027194$ $\approx\qty{368}{m}$}
\bkproblem 工場で、機械を動かしているときは88\,dB、止めているとき（背景の音だけ）は84\,dB でした。機械だけの音は何 dB ですか（⑦・\bkreidairef{9}）。
\bkpans{強さの比で $10^{8.8}-10^{8.4}$ $\approx6.31\times10^8-2.51\times10^8$ $=3.80\times10^8$。$10\log_{10}(3.80\times10^8)$ $\approx\qty{85.8}{dB}$}
\bkproblem \bkreidairef{14}の教室（\qty{189}{m^3}、吸音力\qty{30}{m^2}）の残響時間を0.6秒にしたい。吸音力をあと何 m$^2$ ふやせばよいですか。天井（\qty{63}{m^2}、今の吸音率0.02）を吸音板にはりかえるとき、吸音板の吸音率はいくら要りますか（\nobreak\bkreidairef{14}\nobreak）。
\bkpans{要る吸音力 $0.161\times189\div0.6\approx\qty{50.7}{m^2}$。あと約\qty{20.7}{m^2}。$0.02+20.7\div63\approx0.35$ 以上}
\bkproblem 音圧\qty{1}{Pa}・\qty{100}{Hz}の音では、空気のずれの振幅（実効値の $\sqrt2$ 倍）は何 mm ですか。同じ音圧の\qty{1000}{Hz}の音と比べましょう（\nobreak\bkreidairef{3}\nobreak）。
\bkpans{$u=1\div413\approx\qty{2.42e-3}{m/s}$。振幅 $\sqrt2\times2.42\times10^{-3}$ $\div(2\pi\times100)$ $\approx\qty{5.4e-6}{m}$（約\qty{0.0054}{mm}）。\qty{1000}{Hz}の10倍}
\bkproblem まっすぐな道路（線の音源）から\qty{10}{m}の家で、外の音が\qty{500}{Hz}で75\,dB です。（1）厚さ\qty{5}{mm}のガラス（34\,dB）の窓を、厚さ\qty{10}{mm}のガラスにかえると、部屋の中は何 dB になりますか（質量則）。（2）窓は\qty{5}{mm}のまま、道路から\qty{40}{m}の家に住むとどうですか。（3）両方を合わせるとどうですか（⑦⑨・\bkreidairef{11}・\bkreidairef{13}）。
\bkpans{（1）重さ2倍で $+6$\,dB、透過損失40\,dB なので、$75-40=\qty{35}{dB}$（\qty{5}{mm}のままなら41\,dB）。（2）距離4倍で外は $75-6=\qty{69}{dB}$、中は $69-34=\qty{35}{dB}$。窓を2倍の重さにするのと、4倍離れるのが同じくらい効く。（3）$69-40=\qty{29}{dB}$（dB の引き算が2回重なる）}

%% 修3-4：ふりかえりと吹き出しを同じページに組むため
\Needspace{16\baselineskip}
\begin{furikaeri}
\begin{checklist}
\item 雷やこだまの時間から、音の速さを使って距離を求められる。\bkfurref{ph:波-9:sec2}
\item 音の大きさを dB で表し、音を合わせたときや、離れたときの大きさを計算できる。\bkfurref{ph:波-9:sec7}
\item 縮みにくさと密度から気体・液体・固体の音の速さを求め、遮音と残響を見積もれる。\bkfurref{ph:波-9:sec3}
\end{checklist}
\end{furikaeri}
\bkfukidashi{mina}{音の速さと大きさを、数ではかれるようになったね。次の波-10 では、耳に聞こえない超音波のこだまで、体の中や鉄の中を見るしくみを調べるよ。}
